% =====================================================================
%  TRB 2027 · MAIN PAPER SKELETON (draft)
%  连续动作空间下可证明 COLREGs【方向】合规的无人船避碰 · 投影式安全盾 + RL
%
%  这是把论文从模板占位文(0626-paper/paper.tex = Lorem ipsum)推进成【真骨架】的草稿。
%  每节：已定稿处写草稿正文；待实验数据处标 [PENDING: ...]；诚实边界处标 [HONESTY]。
%  可证明层正文由独立文件提供：% =====================================================================
%  TRB 2027 · Paper section draft: THE PROVABLE SAFETY LAYER
%  论文正式章节草稿（可证明层）· 由三份理论底稿忠实落稿：
%    - Paper/可证明层_正式命题_0707.md      (单障碍 命题 1/2/3)
%    - Paper/可证明层_多障碍扩展_0716.md    (多障碍 ∀i / ∃i 扩展)
%    - Paper/失败机理与紧急兜底刻画_0714.md (Limitations / 紧急兜底)
%  溯源：03 决策日志 L133-L141（原推导）· L164（诚实审计）· L165（gap#1 攻坚）·
%        L166（证书接进 repo）· L189 H（多障碍）· L192（承重订正）。
%
%  诚实红线（继承 CLAUDE.md §0 + 三份底稿）——本节任何改写都必须守：
%    绝不写 provably collision-free / provably compliant（裸）/ complete provable coverage。
%    能 claim 的只有：远场单步无碰(∀i·sound) + 每让路步方向合规(by-construction·作用域受限)
%    + 迫近不可避 sound 证书(∃i)；恒速(CV)单/多障碍模型下基准经验零船体碰撞。可证明档 = 中。
%
%  用法：% =====================================================================
%  TRB 2027 · Paper section draft: THE PROVABLE SAFETY LAYER
%  论文正式章节草稿（可证明层）· 由三份理论底稿忠实落稿：
%    - Paper/可证明层_正式命题_0707.md      (单障碍 命题 1/2/3)
%    - Paper/可证明层_多障碍扩展_0716.md    (多障碍 ∀i / ∃i 扩展)
%    - Paper/失败机理与紧急兜底刻画_0714.md (Limitations / 紧急兜底)
%  溯源：03 决策日志 L133-L141（原推导）· L164（诚实审计）· L165（gap#1 攻坚）·
%        L166（证书接进 repo）· L189 H（多障碍）· L192（承重订正）。
%
%  诚实红线（继承 CLAUDE.md §0 + 三份底稿）——本节任何改写都必须守：
%    绝不写 provably collision-free / provably compliant（裸）/ complete provable coverage。
%    能 claim 的只有：远场单步无碰(∀i·sound) + 每让路步方向合规(by-construction·作用域受限)
%    + 迫近不可避 sound 证书(∃i)；恒速(CV)单/多障碍模型下基准经验零船体碰撞。可证明档 = 中。
%
%  用法：% =====================================================================
%  TRB 2027 · Paper section draft: THE PROVABLE SAFETY LAYER
%  论文正式章节草稿（可证明层）· 由三份理论底稿忠实落稿：
%    - Paper/可证明层_正式命题_0707.md      (单障碍 命题 1/2/3)
%    - Paper/可证明层_多障碍扩展_0716.md    (多障碍 ∀i / ∃i 扩展)
%    - Paper/失败机理与紧急兜底刻画_0714.md (Limitations / 紧急兜底)
%  溯源：03 决策日志 L133-L141（原推导）· L164（诚实审计）· L165（gap#1 攻坚）·
%        L166（证书接进 repo）· L189 H（多障碍）· L192（承重订正）。
%
%  诚实红线（继承 CLAUDE.md §0 + 三份底稿）——本节任何改写都必须守：
%    绝不写 provably collision-free / provably compliant（裸）/ complete provable coverage。
%    能 claim 的只有：远场单步无碰(∀i·sound) + 每让路步方向合规(by-construction·作用域受限)
%    + 迫近不可避 sound 证书(∃i)；恒速(CV)单/多障碍模型下基准经验零船体碰撞。可证明档 = 中。
%
%  用法：\input{sec_provable_safety_layer} 进 Methods 之后（或作为独立 section）。
%  依赖宏包：\usepackage{amsmath,amsthm,amssymb,bm}。TRB 无附录/7500 词 →
%  正文留命题陈述 + 证明梗概 + 诚实作用域；完整证明移交外部技术报告（见 \cite{techreport}）。
% =====================================================================

% ---- 若主文档尚未定义，取消注释以下 theorem 环境 ----
% \newtheorem{proposition}{Proposition}
% \newtheorem{assumption}{Assumption}
% \theoremstyle{remark}\newtheorem*{scope}{Scope and honest limits}

\section{The Provable Safety Layer}
\label{sec:provable}

Our contribution is \emph{not} a claim of provable collision freedom. Prior
provably COLREGs-\emph{direction}-compliant collision avoidance either sacrificed
the continuous action space (discrete action masking, Krasowski and Althoff 2024)
or sacrificed the guarantee (continuous-action work reverts to soft reward
shaping). We place a \emph{projection-based safety shield} on maritime COLREGs in
the continuous action space and characterize \emph{exactly what can and cannot be
proven}. The safety layer is organized by encounter regime and yields three
guarantees of different strengths: a far-field single-step no-collision guarantee
(Proposition~\ref{prop:farfield}), a per-give-way-step direction-compliance
guarantee (Proposition~\ref{prop:compliance}), and a \emph{sound} imminent-%
unavoidability certificate for the emergency regime (Proposition~\ref{prop:imminent}).
Each is stated with its full assumption set, and each carries an explicit scope
boundary. Full proofs are deferred to the external technical report; here we give
statements and proof sketches. Throughout, ``proven'' means a strict lemma;
existence-level and empirical results are labeled as such.

\subsection{System setting and notation}
\label{sec:provable-setup}

We adopt the yaw-constrained point-mass vessel dynamics of Krasowski and Althoff
(2024) (vessel class SR108 container ship, $175 \times 25.4$~m):
\begin{equation}
  \dot{x}=v\cos\theta,\quad \dot{y}=v\sin\theta,\quad \dot{\theta}=\omega,\quad \dot{v}=a,
  \label{eq:dyn}
\end{equation}
with control bounds $|a|\le a_{\max}=0.24~\mathrm{m/s^2}$,
$|\omega|\le\omega_{\max}=0.03~\mathrm{rad/s}$, speed clipped by the environment to
$0\le v\le v_{\max}=9.5~\mathrm{m/s}$ each step (within a step $v$ may briefly reach
$v_{\max}+a_{\max}\Delta t=11.9~\mathrm{m/s}$), and decision step
$\Delta t=10~\mathrm{s}$. The ego body is the $175\times25.4$~m rectangle with
circumscribed radius $R_{\mathrm{circ}}=\tfrac12\sqrt{175^2+25.4^2}=88.4$~m and
inscribed radius $r_{\mathrm{insc}}=\tfrac12\cdot25.4=12.7$~m. Target vessels occupy
their body rectangle (or, conservatively, their circumscribed disk); unless stated
otherwise, in stand-on regimes they are extrapolated at \emph{constant velocity}
(CV). We write $p_e,p_m$ for ego/target positions,
$\bm{v}_e=v_e(\cos\theta_e,\sin\theta_e)$ for the ego velocity vector, and
$\hat{h},\hat{h}^{\perp}$ for the heading unit vector and its normal.

\subsection{Far-field regime: single-step no-collision}
\label{sec:provable-farfield}

\begin{proposition}[Far-field single-step no-collision]
\label{prop:farfield}
Let $d=\lVert p_e-p_m\rVert$ be the ego--target center distance for a single target.
If
\begin{equation}
  d \;\ge\; d_{\mathrm{safe}} + v_{\mathrm{obs,max}}\,\Delta t + \rho_{\mathrm{ego}},
  \label{eq:farfield}
\end{equation}
then after one decision step the two (circumscribed-disk) bodies cannot intersect,
i.e.\ a collision is impossible in that step. Here
$d_{\mathrm{safe}}=R_{\mathrm{circ,ego}}+R_{\mathrm{circ,obs}}$ and
$\rho_{\mathrm{ego}}=v_{\max}\Delta t+\tfrac12 a_{\max}\Delta t^2$ is the maximum
single-step ego displacement.
\end{proposition}

Substituting constants gives $R_{\mathrm{circ,ego}}=88.4$~m and, because the target
width is not in the state, a conservative square upper bound with the Table~II
safety margin $d_{\mathrm{obs,safety}}=2\ell$:
$R_{\mathrm{circ,obs}}=\tfrac12\sqrt2\,\ell+d_{\mathrm{obs,safety}}
=123.74+350=473.74$~m, so $d_{\mathrm{safe}}=562.16$~m,
$\rho_{\mathrm{ego}}=9.5\cdot10+\tfrac12\cdot0.24\cdot100=107.0$~m, and
$v_{\mathrm{obs,max}}\Delta t=95$~m, yielding a threshold of $764.16$~m
(engineering value $\sim\!780$~m with margin).

\emph{Assumptions.} (A1) single target (for $N$ targets, each $i$ must satisfy the
threshold, Proposition~\ref{prop:farfield-N}); (A2) target speed
$v_{\mathrm{obs}}\le9.5~\mathrm{m/s}$---Krasowski sets $v_{\max}$ only for the ego, so
this bound must be assumed explicitly. It is \emph{empirically verified} on the full
benchmark: across all $2000$ HandcraftedTwoVessel scenarios ($342{,}000$ obstacle
states) the maximum is $7.10~\mathrm{m/s}$, with zero states exceeding $9.5$
($\sim\!2.4~\mathrm{m/s}$ margin); higher-speed real AIS data would require raising
the threshold. (A3) $v\le v_{\max}$ at step end (environment clip); (A4) the target
is approximately CV within a step (stand-on); (A5) circumscribed-disk occupancy
(disks disjoint $\Rightarrow$ bodies disjoint).

\emph{Proof sketch.} By the triangle inequality the next-step center distance
satisfies $d_{\mathrm{next}}\ge d-\lVert\Delta p_e\rVert-\lVert\Delta p_m\rVert
\ge d-\rho_{\mathrm{ego}}-v_{\mathrm{obs,max}}\Delta t$. With \eqref{eq:farfield},
$d_{\mathrm{next}}\ge d_{\mathrm{safe}}=R_{\mathrm{circ,ego}}+R_{\mathrm{circ,obs}}$,
so the circumscribed disks---hence the bodies---are disjoint. $\square$

\begin{scope}
The triangle inequality is direction-independent, so head-on is already the true
worst case and there is no geometric gap; this is the \emph{soundest} block of the
provable layer (an elementary lemma). It is a conservative single-step
\emph{sufficient} condition---not invariance, not a near-field guarantee---and is
currently not implemented as an early exit (an $O(1)$ far-field fast path could be
added).
\end{scope}

\subsection{Give-way regime: per-step direction compliance}
\label{sec:provable-compliance}

\begin{proposition}[Per-give-way-step direction compliance]
\label{prop:compliance}
At every decision step that the state machine classifies as a give-way regime
($\rho\in\{\text{head-on},\text{crossing},\text{overtake}\}$) and for which the
projection quadratic program (QP) is feasible (\texttt{needs\_fallback}=False), the
executed safe action $u_{\mathrm{safe}}$ lies in the COLREGs-mandated compliant
half-plane (e.g.\ give-way starboard turn $\omega\le-\omega_{\mathrm{turn}}$).
\end{proposition}

\emph{Assumptions.} (B1) single target; (B2) correct state-machine classification of
$\rho$ (inheriting the $\sim\!12.6\%$ classification gap and the Lemma-1 tie-break
issue); (B3) QP feasibility (not fallback); (B4) scope $=$ give-way encounters. The
empirical give-way trigger rate is $\sim\!13$--$16\%$ \emph{at the
encounter/episode level} (measured twice over $933$ scenarios; overtake $=0$)---not
per step; stand-on violations are counted per step.

\emph{Proof.} The compliant half-plane enters the QP as a hard interval constraint
($w_{\mathrm{hi}}=-\omega_{\mathrm{turn}}$, etc.), so any projected solution
satisfies it. $\square$ (by construction; no separate theorem required).

\begin{scope}
This does \emph{not} cover the emergency regime $\rho_5$ (which occupies
$\sim\!55$--$60\%$ of conflict steps and whose \texttt{is\_emergency} predicate
structurally bypasses the compliance constraint), nor the relaxed/%
collision-minimizing fallbacks (which drop the half-plane). We therefore write only
``provably give-way-direction-compliant, per give-way step'' and \textbf{never}
``provably compliant / compliant at every step'' (the measured violation rate on the
golden 10-seed run is a nonzero $1.13$--$1.8$ per episode).
\end{scope}

\subsection{Emergency regime: a sound imminent-unavoidability certificate}
\label{sec:provable-imminent}

The emergency controller is an empirical last resort. To characterize it honestly we
do not claim it avoids collisions; instead we provide a \emph{sound sufficient
condition} that certifies when a collision is physically unavoidable, so that the
``black box'' becomes a decidable boundary. This replaces an earlier
turn-only condition that was \emph{unsound} because it ignored acceleration/%
deceleration escapes (gap\#1).

\begin{proposition}[Imminent-unavoidability certificate, single CV target]
\label{prop:imminent}
Consider a single constant-velocity (stand-on) target. Define the ego
reachable-center over-approximation $R_{\mathrm{box}}(t)$ (a heading-frame rectangle
centered at the CV extrapolation $c_e(t)=p_e+\bm{v}_e t$) with longitudinal
half-axis $B_{\mathrm{lon}}(t)=\tfrac12 a_{\mathrm{eff}}t^2$ and lateral half-axis
$L_{\mathrm{lat}}(t)$, where
\begin{align}
  a_{\mathrm{eff}} &= \sqrt{a_{\max}^2+(v_{\mathrm{bnd}}\,\omega_{\max})^2},
    \quad v_{\mathrm{bnd}}=v_{\max}+a_{\max}\Delta t=11.9~\mathrm{m/s},\\
  L_{\mathrm{lat}}(t) &=
    \begin{cases}
      \dfrac{v_{\mathrm{bnd}}}{\omega_{\max}}\bigl(1-\cos\omega_{\max}t\bigr),
        & \omega_{\max}t\le\tfrac{\pi}{2},\\[1.2ex]
      \dfrac{v_{\mathrm{bnd}}}{\omega_{\max}}
        + v_{\mathrm{bnd}}\Bigl(t-\tfrac{\pi}{2\omega_{\max}}\Bigr),
        & \text{otherwise.}
    \end{cases}
\end{align}
Let $O(t)$ be the target body rectangle (center $p_m+\bm{v}_m t$, heading
$\theta_m$). If
\begin{equation}
  \exists\, t^\ast\in[0,T]:\quad
  R_{\mathrm{box}}(t^\ast)\;\subseteq\;\bigl(O(t^\ast)\oplus \mathrm{disk}(r_{\mathrm{insc,ego}}=12.7~\mathrm{m})\bigr),
  \label{eq:imminent}
\end{equation}
then the collision is \emph{unavoidable}: every feasible control sequence collides
with the target at $t^\ast$.
\end{proposition}

\emph{Assumptions.} (A1) single target, \emph{per-pair} (per-pair unavoidable
$\Rightarrow$ global unavoidable; per-pair \emph{avoidable} does \emph{not}
compose); (A2) target CV stand-on (a maneuvering target changes $O(t)$ and voids the
guarantee---we state ``unavoidable within the CV world''); (A3) $v\le
v_{\mathrm{bnd}}=11.9$ within a step (includes the clip overshoot); (A4) the ego
inscribed disk $r_{\mathrm{insc}}=12.7~\mathrm{m}\subseteq$ body at any heading
(from the short edge; the long-edge zero-margin case is pending an independent re-check); (A5)
the certificate fires only in the region $\omega_{\max}t^\ast\le\pi/2$
($\le52.4$~s)---we verified that firing is limited by the longitudinal quadratic term
to $t_{\mathrm{crit,lon}}=\sqrt{2(\tfrac12\cdot175+12.7)/a_{\mathrm{eff}}}=21.6~\mathrm{s}<52.4~\mathrm{s}$;
(A6) single-instant containment (sufficient, not necessary); (A7) occupancy
under-approximation: $\text{obs\_width}\le$ true target width and $L_{\mathrm{obs}}\le$
true target length (a soundness precondition the caller must guarantee---passing
sizes larger than the truth over-approximates $O$ and can yield false positives).

\emph{Proof sketch (three steps).} (i) $R_{\mathrm{box}}$ over-approximates the true
reachable center: $\lVert d\bm{v}/dt\rVert=\sqrt{a^2+(v\omega)^2}\le a_{\mathrm{eff}}$,
and double integration bounds the longitudinal offset by $\tfrac12 a_{\mathrm{eff}}t^2$
while $|y'(t)|\le\int_0^t v_{\mathrm{bnd}}\sin(\omega_{\max}s)\,ds=L_{\mathrm{lat}}(t)$.
(ii) The inscribed disk under-approximates the body. (iii) If
$R_{\mathrm{box}}(t^\ast)\subseteq O(t^\ast)\oplus\mathrm{disk}(12.7)$, then every
reachable ego center lies within $12.7$~m of $O(t^\ast)$, so the ego inscribed disk
($\subseteq$ body) intersects $O(t^\ast)$: a collision. Since this holds for
\emph{all} reachable positions, every trajectory collides at $t^\ast$. $\square$

\emph{Verification.} A six-agent adversarial audit ($20{,}000$ random $+$ $735$
curated cases, with a triple-switch bang-bang escape search under exact hull-collision
geometry) found \emph{zero} true false positives; after hardening the
$v_{\mathrm{bnd}}/L_{\mathrm{lat}}$ envelope, a battery plus $1500$ fuzz cases
remained sound (0 false positives) and non-vacuous (firing at $t^\ast\le7$~s).

\begin{scope}
The certificate is a conservative \emph{one-sided sufficient} condition for
imminent unavoidability. It is \textbf{not} an early-warning cost (it is a terminal
classifier, not a planner), \textbf{not} complete or necessary (no-fire $\ne$
avoidable), and by itself certifies nothing about collision freedom. We
\textbf{never} write ``provably zero-collision'': this is an emergency-regime
classifier, not a system-wide guarantee.
\end{scope}

\subsection{Forward-invariant recursive feasibility: the provable upgrade}
\label{sec:provable-forward-invariant}

Propositions~\ref{prop:farfield}--\ref{prop:compliance} are \emph{one-step}
guarantees; a single-step no-collision constraint is not by itself recursively
feasible (a state may be safe now yet admit no safe continuation). This subsection
upgrades the shield from one-step lookahead to a \emph{forward-invariant} guarantee
on a certified clearable set, which is the paper's central provable contribution and
the differentiator against reactive CBF-QP filters (which, under a $\Delta=10$\,s
zero-order hold, carry no recursive-feasibility proof).

\paragraph{Clearable set.} Fix a single constant-velocity (CV) target. Call a finite
control sequence $m$ a \emph{certified straight-tail escape} for state $s$ if its
tail segment is \emph{constant velocity and straight} ($\omega=0\wedge a=0$) and the
sound Lipschitz clearance certificate (the same interval lower bound used for
Prop.~\ref{prop:farfield}, Section~\ref{sec:provable-farfield}) shows hull
distance $d>0$ on $[0,H]$ with a past-closest-point-of-approach (CPA) margin. Define
the \emph{clearable set} $A:=\{s:\exists\,\text{certified straight-tail escape}\}$.

\begin{lemma}[Constant-velocity tail $\Rightarrow$ permanent clearance]
\label{lem:convex-tail}
If the tail is constant velocity and straight and the target is CV, the hull distance
$g(t)=\mathrm{dist}(\text{two fixed-orientation rectangles translating at constant
relative velocity})$ is \emph{convex} in $t$; hence once $g$ is increasing past the
CPA it increases for all later $t$, so $d>0$ on $[0,H]$ together with the past-CPA
margin implies $d>0$ for all $t\ge0$.
\end{lemma}
\emph{Remark (soundness fix).} Convexity requires the tail to be constant
\emph{velocity}: $\omega=0$ alone is insufficient, since an accelerating straight tail
traces a path that is non-affine in $t$ and $g$ need not be convex. The certificate
therefore requires $a=0$ as well. [HONESTY] The certificate is a \emph{sufficient}
condition; it is sound (never certifies a colliding escape---verified by a
long-horizon fuzz with zero false positives) but conservative (it rejects some safe
escapes).

\begin{proposition}[Control invariance of the clearable set]
\label{prop:forward-invariant}
For every $s\in A$ with certified escape $m=(u_0,u_1,\dots)$, applying the first
control $u_0$ over one decision step (official semantics, with the end-of-step speed
clamp aligned to $\Delta$) yields a successor $s'=F(s,u_0)$ whose tail
$m'=(u_1,u_2,\dots)$ is a certified straight-tail escape for $s'$ (it is the second
half of the \emph{same} physical trajectory, which passed one sound certificate).
Hence $s'\in A$: $A$ is control invariant. Consequently a terminal constraint
``successor $\in A$'', realized as an $O(1)$ re-certification of one concrete backup
maneuver (the backup-maneuver terminal set $U_{\mathrm{term}}$), renders the shielded
closed loop forward invariant on $A$, i.e.\ \emph{provably collision-free for all
future time on $A$ under a single CV target}.
\end{proposition}
\emph{Proof sketch.} The tail is certified because it is a sub-arc of an already
certified trajectory (verified numerically to reproduce the original trajectory's
second half exactly under the aligned clamp). Feasibility of $U_{\mathrm{term}}$ is
non-vacuous: $u_0$ itself keeps the backup certified, so the projection QP always
retains at least the backup as a fallback. $\square$

\paragraph{Compliant forward invariance ($A\cap U_{\mathrm{colregs}}$).}
The shield must pick a give-way-\emph{compliant} first step, so the relevant invariant
is $A\cap U_{\mathrm{colregs}}$. Let the compliant backup be ``turn to the mandated
side for $t_1$\,s, then go straight.'' \textbf{Case 1} ($t_1>\Delta$): the tail still
begins with a compliant turn and remains certified by
Prop.~\ref{prop:forward-invariant}, so $s'\in A\cap U_{\mathrm{colregs}}$---this case
is proved. \textbf{Case 2} ($t_1\le\Delta$, the turn completes within one step): the
tail's first step is straight and a fresh compliant escape from $s'$ is required; a
general proof of this case is open. On the shield's real give-way state distribution,
however, Case 2 \emph{never occurs} (the slow yaw rate, $\sim\!1.7^\circ$/s for a
$175$--$260$\,m hull, forces every compliant clearing turn to exceed one decision
step): across all sampled real give-way states with a compliant backup, $100\%$ fall
in Case 1, the tail remains certified and compliant in $100\%$, and the successor
re-admits a fresh compliant certified escape in $100\%$. We therefore state:
$A\cap U_{\mathrm{colregs}}$ is forward invariant via the proved Case 1 whenever the
compliant backup's turn extends beyond one decision step, a condition met on $100\%$
of the observed give-way distribution.

\begin{scope}
[HONESTY] (i) Single CV target (multi-target is weaker, $A=\bigcap_i A_i$, and per-%
target funnels may conflict). (ii) Membership in $A$ is decided by a finite
constant-control maneuver family plus the sound-but-conservative certificate, so $A$
is an \emph{inner} approximation of the true clearable set; the residual is exactly
the $A$-membership boundary---states not clearable in \emph{any} direction---which
falls to the emergency/relaxed fallback, and is \textbf{not} a
compliance-versus-safety conflict. On the real give-way distribution this residual is
$\approx\!0$--$3\%$ (head-on/overtake $0\%$, crossing $\sim\!2.9\%$), and every
in-$A$ give-way state admits a compliant certified escape. [HONESTY] These rates were
computed with a high-accuracy local integrator that matches the official dynamics to
sub-micron; official-pipeline reconfirmation is [PENDING server, evaluation-only].
(iii) The backup-maneuver terminal constraint $U_{\mathrm{term}}$ is implemented and
unit-tested as a sound module; wiring it into the deployed shield and validating the
closed loop is [PENDING server, evaluation-only]. (iv) The give-way states are drawn
from a well-behaved shielded rollout in which holding course never collides within the
horizon, so these rates characterize the clearable-set geometry, not performance under
adversarially imminent conflict, which needs a give-way-rich adversarial source
[PENDING].
\end{scope}

\subsection{Robust extension to multiple targets}
\label{sec:provable-multi}

For $N$ targets with occupancies $O_i(t)$, collision freedom is a conjunction over
targets, $\text{no collision}\iff\bigwedge_{i=1}^{N}(\text{ego\_body}\cap O_i=\varnothing)$.
This structure is the mathematical root of the robust extension: sufficient
conditions compose as $\forall i$, while the unavoidability certificate composes as
$\exists i$.

\begin{proposition}[Far-field, multi-target]
\label{prop:farfield-N}
If $d_i\ge d_{\mathrm{safe}}+v_{\mathrm{obs,max}}\Delta t+\rho_{\mathrm{ego}}$ for
\emph{all} $i$ (the same $\sim\!780$~m threshold), then after one step the ego body
intersects no target body.
\end{proposition}
The single-target proof holds independently for each $i$; the conjunction is exactly
the definition of collision freedom, so the extension is \emph{sound and tight}
(no extra conservatism beyond the per-pair circumscribed-disk bound). The new
obligation is (A2-N): $v_{\mathrm{obs}}^i\le9.5$ for \emph{each} target---the
single-target empirical bound must be re-verified on multi-target scenarios.

\begin{proposition}[Give-way, multi-target]
\label{prop:compliance-N}
Let $G=\{i:\rho_i \text{ is give-way}\}$ with compliant half-planes $H_i$. At every
step where $(i)$ $H_G:=\bigcap_{i\in G}H_i\ne\varnothing$ (the give-way directions
coexist) and $(ii)$ the QP over $U_{\mathrm{box}}\cap H_G\cap\bigcap_i
U_{\mathrm{collision\text{-}free}}^i$ is feasible, the executed $u_{\mathrm{safe}}$ is
direction-compliant for all $i\in G$.
\end{proposition}
The half-planes enter one QP and the projected solution satisfies their
intersection. The honest scope is sharper than in the single-target case: if two
targets demand opposite give-way directions then $H_G=\varnothing$ and the step falls
to fallback (a COLREGs \emph{violation}). This is not a defect but a boundary of
COLREGs itself---a pairwise rule set with no crisp multi-vessel norm (Rule~2 /
special-circumstances); our defensible position is ``compliant when the pairwise
norms coexist, collision-avoidance-first otherwise.'' As $N$ grows the feasible set
shrinks and the give-way compliance rate \emph{decreases}; this must be measured,
not assumed.

\begin{proposition}[Imminent unavoidability, multi-target]
\label{prop:imminent-N}
With $R_{\mathrm{box}}(t)$ as in Proposition~\ref{prop:imminent} (target-independent),
if $\exists\,i,\ \exists\,t^\ast$ ($\omega_{\max}t^\ast\le\pi/2$) with
$R_{\mathrm{box}}(t^\ast)\subseteq O_i(t^\ast)\oplus\mathrm{disk}(12.7)$, then the
collision is unavoidable.
\end{proposition}
The per-pair certificate for target $i$ implies every reachable trajectory collides
with $O_i$, hence with some target; the existential $\exists i$ is exactly what the
emergency layer needs, at cost $O(N\cdot|T|)$. The half that does \emph{not} compose
is avoidability: per-pair no-fire for all $i$ does \emph{not} imply global
avoidability (evading $A$ may force the ego into the future position of $B$), so
in the multi-target world no-fire is even weaker than for a single target.

\subsection{Honest positioning of the provable layer}
\label{sec:provable-summary}

Table~\ref{tab:provable} summarizes the guarantees, their status, and their scope.
The headline claim we can defend, with all assumptions attached, is: continuous
action $\cap$ provably \emph{direction}-compliant $\cap$ COLREGs (a previously empty
intersection), plus far-field single-step provable safety ($\forall i$, sound), a
sound imminent-unavoidability certificate for the emergency regime ($\exists i$), and
---the central upgrade---\emph{forward-invariant recursive feasibility} on a
certificate-defined clearable set (Prop.~\ref{prop:forward-invariant}): the shield is
provably collision-free for all future time on that set under a single CV target, a
guarantee reactive CBF-QP filters lack under a $10$\,s hold. We do \textbf{not} claim
bare provable collision freedom, complete provable coverage, maneuvering targets, or
multi-target direction conflicts (which fall to fallback by the pairwise nature of
COLREGs); the forward-invariance is scoped to the clearable set, whose residual
($\approx\!0$--$3\%$) is deferred to fallback. The provable tier is therefore
\emph{medium-to-strong}: the invariance upgrade and gap\#1 are sound, but coverage is
still an inner approximation, the closed-loop $U_{\mathrm{term}}$ wiring and official
reconfirmation are [PENDING server, eval-only], and the empirical zero-collision
backing remains single-target until multi-vessel benchmark data are available.

\begin{table}[h]
  \caption{Provable safety layer: guarantees, status, and honest scope}
  \label{tab:provable}
  \centering
  \small
  \begin{tabular}{@{}p{0.20\linewidth}p{0.24\linewidth}p{0.46\linewidth}@{}}
    \toprule
    \textbf{Guarantee} & \textbf{Status} & \textbf{Scope / limits} \\
    \midrule
    Far-field single-step no-collision (Prop.~\ref{prop:farfield},
    \ref{prop:farfield-N})
      & Proven (elementary; single-target bound conservative, multi-target conjunction adds no extra slack)
      & $d\ge\sim\!780$~m; CV target(s); $\forall i$ conjunction; each
        $v_{\mathrm{obs}}\le9.5$; single step, not invariance; early exit not yet
        implemented. \\
    \addlinespace
    Per-give-way-step direction compliance (Prop.~\ref{prop:compliance},
    \ref{prop:compliance-N})
      & By construction (scope-limited)
      & Give-way encounters ($\sim\!13$--$16\%$, encounter-level, not per step);
        excludes emergency $\rho_5$ and fallbacks; multi-target rate drops with $N$;
        direction conflict $\to$ fallback (COLREGs pairwise boundary). Never ``every
        step.'' \\
    \addlinespace
    Imminent-unavoidability certificate (gap\#1) (Prop.~\ref{prop:imminent},
    \ref{prop:imminent-N})
      & Sound sufficient condition ($>\!20{,}000$ adversarial scenarios, 0 false positives)
      & Imminent only; CV target(s); $\exists i$ composes, avoidability does not;
        not complete / not early-warning; a classifier, not a system guarantee. \\
    \addlinespace
    Forward-invariant recursive feasibility on the clearable set $A$
    (Lem.~\ref{lem:convex-tail}, Prop.~\ref{prop:forward-invariant})
      & Proven (control invariance of $A$; certificate sound, 0 false positives on
        long-horizon fuzz). $U_{\mathrm{term}}$ module unit-tested; closed-loop wiring
        [PENDING server, eval-only]
      & Single CV target; $A$ = inner approximation (finite maneuver family +
        conservative certificate); residual $\approx\!0$--$3\%$ = states clearable in
        no direction $\to$ fallback (not a compliance conflict); $A\cap U_{\mathrm{colregs}}$
        invariance proved for Case 1 and holds on $100\%$ of the observed give-way
        distribution; numbers local-integrator, official reconfirmation [PENDING]. \\
    \bottomrule
  \end{tabular}
\end{table}

\paragraph{Emergency fallback and the failure-mechanism boundary.}
When the projection is infeasible ($P=\varnothing$) in the emergency regime
$\rho_5$, the \texttt{is\_emergency} predicate structurally bypasses the direction
constraint and Algorithm~1 minimizes collision risk---an \emph{empirical} fallback,
not a guarantee. On an adversarial purely-geometric baseline that forces the shield
into the emergency branch on $44$--$69\%$ of steps, we observe $45\%$ arrival and
two collisions; auditing confirms both occur with
\texttt{source=emergency},~$\rho=5$, i.e.\ the emergency last resort trying and
failing to avoid, and both lie outside the domain of Proposition~\ref{prop:farfield}
(far-field, non-emergency, single-step). They are not counterexamples to the headline
but they do substantiate that the emergency regime is not a guarantee. On the learned
policy (non-adversarial), we observe zero hull collisions over $10$ seeds
$\times$ $40$ episodes ($400$ total), so the ``empirically zero-collision'' statement holds; we
nonetheless \textbf{never} write provably collision-free and label the emergency
regime as an empirical fallback throughout.
 进 Methods 之后（或作为独立 section）。
%  依赖宏包：\usepackage{amsmath,amsthm,amssymb,bm}。TRB 无附录/7500 词 →
%  正文留命题陈述 + 证明梗概 + 诚实作用域；完整证明移交外部技术报告（见 \cite{techreport}）。
% =====================================================================

% ---- 若主文档尚未定义，取消注释以下 theorem 环境 ----
% \newtheorem{proposition}{Proposition}
% \newtheorem{assumption}{Assumption}
% \theoremstyle{remark}\newtheorem*{scope}{Scope and honest limits}

\section{The Provable Safety Layer}
\label{sec:provable}

Our contribution is \emph{not} a claim of provable collision freedom. Prior
provably COLREGs-\emph{direction}-compliant collision avoidance either sacrificed
the continuous action space (discrete action masking, Krasowski and Althoff 2024)
or sacrificed the guarantee (continuous-action work reverts to soft reward
shaping). We place a \emph{projection-based safety shield} on maritime COLREGs in
the continuous action space and characterize \emph{exactly what can and cannot be
proven}. The safety layer is organized by encounter regime and yields three
guarantees of different strengths: a far-field single-step no-collision guarantee
(Proposition~\ref{prop:farfield}), a per-give-way-step direction-compliance
guarantee (Proposition~\ref{prop:compliance}), and a \emph{sound} imminent-%
unavoidability certificate for the emergency regime (Proposition~\ref{prop:imminent}).
Each is stated with its full assumption set, and each carries an explicit scope
boundary. Full proofs are deferred to the external technical report; here we give
statements and proof sketches. Throughout, ``proven'' means a strict lemma;
existence-level and empirical results are labeled as such.

\subsection{System setting and notation}
\label{sec:provable-setup}

We adopt the yaw-constrained point-mass vessel dynamics of Krasowski and Althoff
(2024) (vessel class SR108 container ship, $175 \times 25.4$~m):
\begin{equation}
  \dot{x}=v\cos\theta,\quad \dot{y}=v\sin\theta,\quad \dot{\theta}=\omega,\quad \dot{v}=a,
  \label{eq:dyn}
\end{equation}
with control bounds $|a|\le a_{\max}=0.24~\mathrm{m/s^2}$,
$|\omega|\le\omega_{\max}=0.03~\mathrm{rad/s}$, speed clipped by the environment to
$0\le v\le v_{\max}=9.5~\mathrm{m/s}$ each step (within a step $v$ may briefly reach
$v_{\max}+a_{\max}\Delta t=11.9~\mathrm{m/s}$), and decision step
$\Delta t=10~\mathrm{s}$. The ego body is the $175\times25.4$~m rectangle with
circumscribed radius $R_{\mathrm{circ}}=\tfrac12\sqrt{175^2+25.4^2}=88.4$~m and
inscribed radius $r_{\mathrm{insc}}=\tfrac12\cdot25.4=12.7$~m. Target vessels occupy
their body rectangle (or, conservatively, their circumscribed disk); unless stated
otherwise, in stand-on regimes they are extrapolated at \emph{constant velocity}
(CV). We write $p_e,p_m$ for ego/target positions,
$\bm{v}_e=v_e(\cos\theta_e,\sin\theta_e)$ for the ego velocity vector, and
$\hat{h},\hat{h}^{\perp}$ for the heading unit vector and its normal.

\subsection{Far-field regime: single-step no-collision}
\label{sec:provable-farfield}

\begin{proposition}[Far-field single-step no-collision]
\label{prop:farfield}
Let $d=\lVert p_e-p_m\rVert$ be the ego--target center distance for a single target.
If
\begin{equation}
  d \;\ge\; d_{\mathrm{safe}} + v_{\mathrm{obs,max}}\,\Delta t + \rho_{\mathrm{ego}},
  \label{eq:farfield}
\end{equation}
then after one decision step the two (circumscribed-disk) bodies cannot intersect,
i.e.\ a collision is impossible in that step. Here
$d_{\mathrm{safe}}=R_{\mathrm{circ,ego}}+R_{\mathrm{circ,obs}}$ and
$\rho_{\mathrm{ego}}=v_{\max}\Delta t+\tfrac12 a_{\max}\Delta t^2$ is the maximum
single-step ego displacement.
\end{proposition}

Substituting constants gives $R_{\mathrm{circ,ego}}=88.4$~m and, because the target
width is not in the state, a conservative square upper bound with the Table~II
safety margin $d_{\mathrm{obs,safety}}=2\ell$:
$R_{\mathrm{circ,obs}}=\tfrac12\sqrt2\,\ell+d_{\mathrm{obs,safety}}
=123.74+350=473.74$~m, so $d_{\mathrm{safe}}=562.16$~m,
$\rho_{\mathrm{ego}}=9.5\cdot10+\tfrac12\cdot0.24\cdot100=107.0$~m, and
$v_{\mathrm{obs,max}}\Delta t=95$~m, yielding a threshold of $764.16$~m
(engineering value $\sim\!780$~m with margin).

\emph{Assumptions.} (A1) single target (for $N$ targets, each $i$ must satisfy the
threshold, Proposition~\ref{prop:farfield-N}); (A2) target speed
$v_{\mathrm{obs}}\le9.5~\mathrm{m/s}$---Krasowski sets $v_{\max}$ only for the ego, so
this bound must be assumed explicitly. It is \emph{empirically verified} on the full
benchmark: across all $2000$ HandcraftedTwoVessel scenarios ($342{,}000$ obstacle
states) the maximum is $7.10~\mathrm{m/s}$, with zero states exceeding $9.5$
($\sim\!2.4~\mathrm{m/s}$ margin); higher-speed real AIS data would require raising
the threshold. (A3) $v\le v_{\max}$ at step end (environment clip); (A4) the target
is approximately CV within a step (stand-on); (A5) circumscribed-disk occupancy
(disks disjoint $\Rightarrow$ bodies disjoint).

\emph{Proof sketch.} By the triangle inequality the next-step center distance
satisfies $d_{\mathrm{next}}\ge d-\lVert\Delta p_e\rVert-\lVert\Delta p_m\rVert
\ge d-\rho_{\mathrm{ego}}-v_{\mathrm{obs,max}}\Delta t$. With \eqref{eq:farfield},
$d_{\mathrm{next}}\ge d_{\mathrm{safe}}=R_{\mathrm{circ,ego}}+R_{\mathrm{circ,obs}}$,
so the circumscribed disks---hence the bodies---are disjoint. $\square$

\begin{scope}
The triangle inequality is direction-independent, so head-on is already the true
worst case and there is no geometric gap; this is the \emph{soundest} block of the
provable layer (an elementary lemma). It is a conservative single-step
\emph{sufficient} condition---not invariance, not a near-field guarantee---and is
currently not implemented as an early exit (an $O(1)$ far-field fast path could be
added).
\end{scope}

\subsection{Give-way regime: per-step direction compliance}
\label{sec:provable-compliance}

\begin{proposition}[Per-give-way-step direction compliance]
\label{prop:compliance}
At every decision step that the state machine classifies as a give-way regime
($\rho\in\{\text{head-on},\text{crossing},\text{overtake}\}$) and for which the
projection quadratic program (QP) is feasible (\texttt{needs\_fallback}=False), the
executed safe action $u_{\mathrm{safe}}$ lies in the COLREGs-mandated compliant
half-plane (e.g.\ give-way starboard turn $\omega\le-\omega_{\mathrm{turn}}$).
\end{proposition}

\emph{Assumptions.} (B1) single target; (B2) correct state-machine classification of
$\rho$ (inheriting the $\sim\!12.6\%$ classification gap and the Lemma-1 tie-break
issue); (B3) QP feasibility (not fallback); (B4) scope $=$ give-way encounters. The
empirical give-way trigger rate is $\sim\!13$--$16\%$ \emph{at the
encounter/episode level} (measured twice over $933$ scenarios; overtake $=0$)---not
per step; stand-on violations are counted per step.

\emph{Proof.} The compliant half-plane enters the QP as a hard interval constraint
($w_{\mathrm{hi}}=-\omega_{\mathrm{turn}}$, etc.), so any projected solution
satisfies it. $\square$ (by construction; no separate theorem required).

\begin{scope}
This does \emph{not} cover the emergency regime $\rho_5$ (which occupies
$\sim\!55$--$60\%$ of conflict steps and whose \texttt{is\_emergency} predicate
structurally bypasses the compliance constraint), nor the relaxed/%
collision-minimizing fallbacks (which drop the half-plane). We therefore write only
``provably give-way-direction-compliant, per give-way step'' and \textbf{never}
``provably compliant / compliant at every step'' (the measured violation rate on the
golden 10-seed run is a nonzero $1.13$--$1.8$ per episode).
\end{scope}

\subsection{Emergency regime: a sound imminent-unavoidability certificate}
\label{sec:provable-imminent}

The emergency controller is an empirical last resort. To characterize it honestly we
do not claim it avoids collisions; instead we provide a \emph{sound sufficient
condition} that certifies when a collision is physically unavoidable, so that the
``black box'' becomes a decidable boundary. This replaces an earlier
turn-only condition that was \emph{unsound} because it ignored acceleration/%
deceleration escapes (gap\#1).

\begin{proposition}[Imminent-unavoidability certificate, single CV target]
\label{prop:imminent}
Consider a single constant-velocity (stand-on) target. Define the ego
reachable-center over-approximation $R_{\mathrm{box}}(t)$ (a heading-frame rectangle
centered at the CV extrapolation $c_e(t)=p_e+\bm{v}_e t$) with longitudinal
half-axis $B_{\mathrm{lon}}(t)=\tfrac12 a_{\mathrm{eff}}t^2$ and lateral half-axis
$L_{\mathrm{lat}}(t)$, where
\begin{align}
  a_{\mathrm{eff}} &= \sqrt{a_{\max}^2+(v_{\mathrm{bnd}}\,\omega_{\max})^2},
    \quad v_{\mathrm{bnd}}=v_{\max}+a_{\max}\Delta t=11.9~\mathrm{m/s},\\
  L_{\mathrm{lat}}(t) &=
    \begin{cases}
      \dfrac{v_{\mathrm{bnd}}}{\omega_{\max}}\bigl(1-\cos\omega_{\max}t\bigr),
        & \omega_{\max}t\le\tfrac{\pi}{2},\\[1.2ex]
      \dfrac{v_{\mathrm{bnd}}}{\omega_{\max}}
        + v_{\mathrm{bnd}}\Bigl(t-\tfrac{\pi}{2\omega_{\max}}\Bigr),
        & \text{otherwise.}
    \end{cases}
\end{align}
Let $O(t)$ be the target body rectangle (center $p_m+\bm{v}_m t$, heading
$\theta_m$). If
\begin{equation}
  \exists\, t^\ast\in[0,T]:\quad
  R_{\mathrm{box}}(t^\ast)\;\subseteq\;\bigl(O(t^\ast)\oplus \mathrm{disk}(r_{\mathrm{insc,ego}}=12.7~\mathrm{m})\bigr),
  \label{eq:imminent}
\end{equation}
then the collision is \emph{unavoidable}: every feasible control sequence collides
with the target at $t^\ast$.
\end{proposition}

\emph{Assumptions.} (A1) single target, \emph{per-pair} (per-pair unavoidable
$\Rightarrow$ global unavoidable; per-pair \emph{avoidable} does \emph{not}
compose); (A2) target CV stand-on (a maneuvering target changes $O(t)$ and voids the
guarantee---we state ``unavoidable within the CV world''); (A3) $v\le
v_{\mathrm{bnd}}=11.9$ within a step (includes the clip overshoot); (A4) the ego
inscribed disk $r_{\mathrm{insc}}=12.7~\mathrm{m}\subseteq$ body at any heading
(from the short edge; the long-edge zero-margin case is pending an independent re-check); (A5)
the certificate fires only in the region $\omega_{\max}t^\ast\le\pi/2$
($\le52.4$~s)---we verified that firing is limited by the longitudinal quadratic term
to $t_{\mathrm{crit,lon}}=\sqrt{2(\tfrac12\cdot175+12.7)/a_{\mathrm{eff}}}=21.6~\mathrm{s}<52.4~\mathrm{s}$;
(A6) single-instant containment (sufficient, not necessary); (A7) occupancy
under-approximation: $\text{obs\_width}\le$ true target width and $L_{\mathrm{obs}}\le$
true target length (a soundness precondition the caller must guarantee---passing
sizes larger than the truth over-approximates $O$ and can yield false positives).

\emph{Proof sketch (three steps).} (i) $R_{\mathrm{box}}$ over-approximates the true
reachable center: $\lVert d\bm{v}/dt\rVert=\sqrt{a^2+(v\omega)^2}\le a_{\mathrm{eff}}$,
and double integration bounds the longitudinal offset by $\tfrac12 a_{\mathrm{eff}}t^2$
while $|y'(t)|\le\int_0^t v_{\mathrm{bnd}}\sin(\omega_{\max}s)\,ds=L_{\mathrm{lat}}(t)$.
(ii) The inscribed disk under-approximates the body. (iii) If
$R_{\mathrm{box}}(t^\ast)\subseteq O(t^\ast)\oplus\mathrm{disk}(12.7)$, then every
reachable ego center lies within $12.7$~m of $O(t^\ast)$, so the ego inscribed disk
($\subseteq$ body) intersects $O(t^\ast)$: a collision. Since this holds for
\emph{all} reachable positions, every trajectory collides at $t^\ast$. $\square$

\emph{Verification.} A six-agent adversarial audit ($20{,}000$ random $+$ $735$
curated cases, with a triple-switch bang-bang escape search under exact hull-collision
geometry) found \emph{zero} true false positives; after hardening the
$v_{\mathrm{bnd}}/L_{\mathrm{lat}}$ envelope, a battery plus $1500$ fuzz cases
remained sound (0 false positives) and non-vacuous (firing at $t^\ast\le7$~s).

\begin{scope}
The certificate is a conservative \emph{one-sided sufficient} condition for
imminent unavoidability. It is \textbf{not} an early-warning cost (it is a terminal
classifier, not a planner), \textbf{not} complete or necessary (no-fire $\ne$
avoidable), and by itself certifies nothing about collision freedom. We
\textbf{never} write ``provably zero-collision'': this is an emergency-regime
classifier, not a system-wide guarantee.
\end{scope}

\subsection{Forward-invariant recursive feasibility: the provable upgrade}
\label{sec:provable-forward-invariant}

Propositions~\ref{prop:farfield}--\ref{prop:compliance} are \emph{one-step}
guarantees; a single-step no-collision constraint is not by itself recursively
feasible (a state may be safe now yet admit no safe continuation). This subsection
upgrades the shield from one-step lookahead to a \emph{forward-invariant} guarantee
on a certified clearable set, which is the paper's central provable contribution and
the differentiator against reactive CBF-QP filters (which, under a $\Delta=10$\,s
zero-order hold, carry no recursive-feasibility proof).

\paragraph{Clearable set.} Fix a single constant-velocity (CV) target. Call a finite
control sequence $m$ a \emph{certified straight-tail escape} for state $s$ if its
tail segment is \emph{constant velocity and straight} ($\omega=0\wedge a=0$) and the
sound Lipschitz clearance certificate (the same interval lower bound used for
Prop.~\ref{prop:farfield}, Section~\ref{sec:provable-farfield}) shows hull
distance $d>0$ on $[0,H]$ with a past-closest-point-of-approach (CPA) margin. Define
the \emph{clearable set} $A:=\{s:\exists\,\text{certified straight-tail escape}\}$.

\begin{lemma}[Constant-velocity tail $\Rightarrow$ permanent clearance]
\label{lem:convex-tail}
If the tail is constant velocity and straight and the target is CV, the hull distance
$g(t)=\mathrm{dist}(\text{two fixed-orientation rectangles translating at constant
relative velocity})$ is \emph{convex} in $t$; hence once $g$ is increasing past the
CPA it increases for all later $t$, so $d>0$ on $[0,H]$ together with the past-CPA
margin implies $d>0$ for all $t\ge0$.
\end{lemma}
\emph{Remark (soundness fix).} Convexity requires the tail to be constant
\emph{velocity}: $\omega=0$ alone is insufficient, since an accelerating straight tail
traces a path that is non-affine in $t$ and $g$ need not be convex. The certificate
therefore requires $a=0$ as well. [HONESTY] The certificate is a \emph{sufficient}
condition; it is sound (never certifies a colliding escape---verified by a
long-horizon fuzz with zero false positives) but conservative (it rejects some safe
escapes).

\begin{proposition}[Control invariance of the clearable set]
\label{prop:forward-invariant}
For every $s\in A$ with certified escape $m=(u_0,u_1,\dots)$, applying the first
control $u_0$ over one decision step (official semantics, with the end-of-step speed
clamp aligned to $\Delta$) yields a successor $s'=F(s,u_0)$ whose tail
$m'=(u_1,u_2,\dots)$ is a certified straight-tail escape for $s'$ (it is the second
half of the \emph{same} physical trajectory, which passed one sound certificate).
Hence $s'\in A$: $A$ is control invariant. Consequently a terminal constraint
``successor $\in A$'', realized as an $O(1)$ re-certification of one concrete backup
maneuver (the backup-maneuver terminal set $U_{\mathrm{term}}$), renders the shielded
closed loop forward invariant on $A$, i.e.\ \emph{provably collision-free for all
future time on $A$ under a single CV target}.
\end{proposition}
\emph{Proof sketch.} The tail is certified because it is a sub-arc of an already
certified trajectory (verified numerically to reproduce the original trajectory's
second half exactly under the aligned clamp). Feasibility of $U_{\mathrm{term}}$ is
non-vacuous: $u_0$ itself keeps the backup certified, so the projection QP always
retains at least the backup as a fallback. $\square$

\paragraph{Compliant forward invariance ($A\cap U_{\mathrm{colregs}}$).}
The shield must pick a give-way-\emph{compliant} first step, so the relevant invariant
is $A\cap U_{\mathrm{colregs}}$. Let the compliant backup be ``turn to the mandated
side for $t_1$\,s, then go straight.'' \textbf{Case 1} ($t_1>\Delta$): the tail still
begins with a compliant turn and remains certified by
Prop.~\ref{prop:forward-invariant}, so $s'\in A\cap U_{\mathrm{colregs}}$---this case
is proved. \textbf{Case 2} ($t_1\le\Delta$, the turn completes within one step): the
tail's first step is straight and a fresh compliant escape from $s'$ is required; a
general proof of this case is open. On the shield's real give-way state distribution,
however, Case 2 \emph{never occurs} (the slow yaw rate, $\sim\!1.7^\circ$/s for a
$175$--$260$\,m hull, forces every compliant clearing turn to exceed one decision
step): across all sampled real give-way states with a compliant backup, $100\%$ fall
in Case 1, the tail remains certified and compliant in $100\%$, and the successor
re-admits a fresh compliant certified escape in $100\%$. We therefore state:
$A\cap U_{\mathrm{colregs}}$ is forward invariant via the proved Case 1 whenever the
compliant backup's turn extends beyond one decision step, a condition met on $100\%$
of the observed give-way distribution.

\begin{scope}
[HONESTY] (i) Single CV target (multi-target is weaker, $A=\bigcap_i A_i$, and per-%
target funnels may conflict). (ii) Membership in $A$ is decided by a finite
constant-control maneuver family plus the sound-but-conservative certificate, so $A$
is an \emph{inner} approximation of the true clearable set; the residual is exactly
the $A$-membership boundary---states not clearable in \emph{any} direction---which
falls to the emergency/relaxed fallback, and is \textbf{not} a
compliance-versus-safety conflict. On the real give-way distribution this residual is
$\approx\!0$--$3\%$ (head-on/overtake $0\%$, crossing $\sim\!2.9\%$), and every
in-$A$ give-way state admits a compliant certified escape. [HONESTY] These rates were
computed with a high-accuracy local integrator that matches the official dynamics to
sub-micron; official-pipeline reconfirmation is [PENDING server, evaluation-only].
(iii) The backup-maneuver terminal constraint $U_{\mathrm{term}}$ is implemented and
unit-tested as a sound module; wiring it into the deployed shield and validating the
closed loop is [PENDING server, evaluation-only]. (iv) The give-way states are drawn
from a well-behaved shielded rollout in which holding course never collides within the
horizon, so these rates characterize the clearable-set geometry, not performance under
adversarially imminent conflict, which needs a give-way-rich adversarial source
[PENDING].
\end{scope}

\subsection{Robust extension to multiple targets}
\label{sec:provable-multi}

For $N$ targets with occupancies $O_i(t)$, collision freedom is a conjunction over
targets, $\text{no collision}\iff\bigwedge_{i=1}^{N}(\text{ego\_body}\cap O_i=\varnothing)$.
This structure is the mathematical root of the robust extension: sufficient
conditions compose as $\forall i$, while the unavoidability certificate composes as
$\exists i$.

\begin{proposition}[Far-field, multi-target]
\label{prop:farfield-N}
If $d_i\ge d_{\mathrm{safe}}+v_{\mathrm{obs,max}}\Delta t+\rho_{\mathrm{ego}}$ for
\emph{all} $i$ (the same $\sim\!780$~m threshold), then after one step the ego body
intersects no target body.
\end{proposition}
The single-target proof holds independently for each $i$; the conjunction is exactly
the definition of collision freedom, so the extension is \emph{sound and tight}
(no extra conservatism beyond the per-pair circumscribed-disk bound). The new
obligation is (A2-N): $v_{\mathrm{obs}}^i\le9.5$ for \emph{each} target---the
single-target empirical bound must be re-verified on multi-target scenarios.

\begin{proposition}[Give-way, multi-target]
\label{prop:compliance-N}
Let $G=\{i:\rho_i \text{ is give-way}\}$ with compliant half-planes $H_i$. At every
step where $(i)$ $H_G:=\bigcap_{i\in G}H_i\ne\varnothing$ (the give-way directions
coexist) and $(ii)$ the QP over $U_{\mathrm{box}}\cap H_G\cap\bigcap_i
U_{\mathrm{collision\text{-}free}}^i$ is feasible, the executed $u_{\mathrm{safe}}$ is
direction-compliant for all $i\in G$.
\end{proposition}
The half-planes enter one QP and the projected solution satisfies their
intersection. The honest scope is sharper than in the single-target case: if two
targets demand opposite give-way directions then $H_G=\varnothing$ and the step falls
to fallback (a COLREGs \emph{violation}). This is not a defect but a boundary of
COLREGs itself---a pairwise rule set with no crisp multi-vessel norm (Rule~2 /
special-circumstances); our defensible position is ``compliant when the pairwise
norms coexist, collision-avoidance-first otherwise.'' As $N$ grows the feasible set
shrinks and the give-way compliance rate \emph{decreases}; this must be measured,
not assumed.

\begin{proposition}[Imminent unavoidability, multi-target]
\label{prop:imminent-N}
With $R_{\mathrm{box}}(t)$ as in Proposition~\ref{prop:imminent} (target-independent),
if $\exists\,i,\ \exists\,t^\ast$ ($\omega_{\max}t^\ast\le\pi/2$) with
$R_{\mathrm{box}}(t^\ast)\subseteq O_i(t^\ast)\oplus\mathrm{disk}(12.7)$, then the
collision is unavoidable.
\end{proposition}
The per-pair certificate for target $i$ implies every reachable trajectory collides
with $O_i$, hence with some target; the existential $\exists i$ is exactly what the
emergency layer needs, at cost $O(N\cdot|T|)$. The half that does \emph{not} compose
is avoidability: per-pair no-fire for all $i$ does \emph{not} imply global
avoidability (evading $A$ may force the ego into the future position of $B$), so
in the multi-target world no-fire is even weaker than for a single target.

\subsection{Honest positioning of the provable layer}
\label{sec:provable-summary}

Table~\ref{tab:provable} summarizes the guarantees, their status, and their scope.
The headline claim we can defend, with all assumptions attached, is: continuous
action $\cap$ provably \emph{direction}-compliant $\cap$ COLREGs (a previously empty
intersection), plus far-field single-step provable safety ($\forall i$, sound), a
sound imminent-unavoidability certificate for the emergency regime ($\exists i$), and
---the central upgrade---\emph{forward-invariant recursive feasibility} on a
certificate-defined clearable set (Prop.~\ref{prop:forward-invariant}): the shield is
provably collision-free for all future time on that set under a single CV target, a
guarantee reactive CBF-QP filters lack under a $10$\,s hold. We do \textbf{not} claim
bare provable collision freedom, complete provable coverage, maneuvering targets, or
multi-target direction conflicts (which fall to fallback by the pairwise nature of
COLREGs); the forward-invariance is scoped to the clearable set, whose residual
($\approx\!0$--$3\%$) is deferred to fallback. The provable tier is therefore
\emph{medium-to-strong}: the invariance upgrade and gap\#1 are sound, but coverage is
still an inner approximation, the closed-loop $U_{\mathrm{term}}$ wiring and official
reconfirmation are [PENDING server, eval-only], and the empirical zero-collision
backing remains single-target until multi-vessel benchmark data are available.

\begin{table}[h]
  \caption{Provable safety layer: guarantees, status, and honest scope}
  \label{tab:provable}
  \centering
  \small
  \begin{tabular}{@{}p{0.20\linewidth}p{0.24\linewidth}p{0.46\linewidth}@{}}
    \toprule
    \textbf{Guarantee} & \textbf{Status} & \textbf{Scope / limits} \\
    \midrule
    Far-field single-step no-collision (Prop.~\ref{prop:farfield},
    \ref{prop:farfield-N})
      & Proven (elementary; single-target bound conservative, multi-target conjunction adds no extra slack)
      & $d\ge\sim\!780$~m; CV target(s); $\forall i$ conjunction; each
        $v_{\mathrm{obs}}\le9.5$; single step, not invariance; early exit not yet
        implemented. \\
    \addlinespace
    Per-give-way-step direction compliance (Prop.~\ref{prop:compliance},
    \ref{prop:compliance-N})
      & By construction (scope-limited)
      & Give-way encounters ($\sim\!13$--$16\%$, encounter-level, not per step);
        excludes emergency $\rho_5$ and fallbacks; multi-target rate drops with $N$;
        direction conflict $\to$ fallback (COLREGs pairwise boundary). Never ``every
        step.'' \\
    \addlinespace
    Imminent-unavoidability certificate (gap\#1) (Prop.~\ref{prop:imminent},
    \ref{prop:imminent-N})
      & Sound sufficient condition ($>\!20{,}000$ adversarial scenarios, 0 false positives)
      & Imminent only; CV target(s); $\exists i$ composes, avoidability does not;
        not complete / not early-warning; a classifier, not a system guarantee. \\
    \addlinespace
    Forward-invariant recursive feasibility on the clearable set $A$
    (Lem.~\ref{lem:convex-tail}, Prop.~\ref{prop:forward-invariant})
      & Proven (control invariance of $A$; certificate sound, 0 false positives on
        long-horizon fuzz). $U_{\mathrm{term}}$ module unit-tested; closed-loop wiring
        [PENDING server, eval-only]
      & Single CV target; $A$ = inner approximation (finite maneuver family +
        conservative certificate); residual $\approx\!0$--$3\%$ = states clearable in
        no direction $\to$ fallback (not a compliance conflict); $A\cap U_{\mathrm{colregs}}$
        invariance proved for Case 1 and holds on $100\%$ of the observed give-way
        distribution; numbers local-integrator, official reconfirmation [PENDING]. \\
    \bottomrule
  \end{tabular}
\end{table}

\paragraph{Emergency fallback and the failure-mechanism boundary.}
When the projection is infeasible ($P=\varnothing$) in the emergency regime
$\rho_5$, the \texttt{is\_emergency} predicate structurally bypasses the direction
constraint and Algorithm~1 minimizes collision risk---an \emph{empirical} fallback,
not a guarantee. On an adversarial purely-geometric baseline that forces the shield
into the emergency branch on $44$--$69\%$ of steps, we observe $45\%$ arrival and
two collisions; auditing confirms both occur with
\texttt{source=emergency},~$\rho=5$, i.e.\ the emergency last resort trying and
failing to avoid, and both lie outside the domain of Proposition~\ref{prop:farfield}
(far-field, non-emergency, single-step). They are not counterexamples to the headline
but they do substantiate that the emergency regime is not a guarantee. On the learned
policy (non-adversarial), we observe zero hull collisions over $10$ seeds
$\times$ $40$ episodes ($400$ total), so the ``empirically zero-collision'' statement holds; we
nonetheless \textbf{never} write provably collision-free and label the emergency
regime as an empirical fallback throughout.
 进 Methods 之后（或作为独立 section）。
%  依赖宏包：\usepackage{amsmath,amsthm,amssymb,bm}。TRB 无附录/7500 词 →
%  正文留命题陈述 + 证明梗概 + 诚实作用域；完整证明移交外部技术报告（见 \cite{techreport}）。
% =====================================================================

% ---- 若主文档尚未定义，取消注释以下 theorem 环境 ----
% \newtheorem{proposition}{Proposition}
% \newtheorem{assumption}{Assumption}
% \theoremstyle{remark}\newtheorem*{scope}{Scope and honest limits}

\section{The Provable Safety Layer}
\label{sec:provable}

Our contribution is \emph{not} a claim of provable collision freedom. Prior
provably COLREGs-\emph{direction}-compliant collision avoidance either sacrificed
the continuous action space (discrete action masking, Krasowski and Althoff 2024)
or sacrificed the guarantee (continuous-action work reverts to soft reward
shaping). We place a \emph{projection-based safety shield} on maritime COLREGs in
the continuous action space and characterize \emph{exactly what can and cannot be
proven}. The safety layer is organized by encounter regime and yields three
guarantees of different strengths: a far-field single-step no-collision guarantee
(Proposition~\ref{prop:farfield}), a per-give-way-step direction-compliance
guarantee (Proposition~\ref{prop:compliance}), and a \emph{sound} imminent-%
unavoidability certificate for the emergency regime (Proposition~\ref{prop:imminent}).
Each is stated with its full assumption set, and each carries an explicit scope
boundary. Full proofs are deferred to the external technical report; here we give
statements and proof sketches. Throughout, ``proven'' means a strict lemma;
existence-level and empirical results are labeled as such.

\subsection{System setting and notation}
\label{sec:provable-setup}

We adopt the yaw-constrained point-mass vessel dynamics of Krasowski and Althoff
(2024) (vessel class SR108 container ship, $175 \times 25.4$~m):
\begin{equation}
  \dot{x}=v\cos\theta,\quad \dot{y}=v\sin\theta,\quad \dot{\theta}=\omega,\quad \dot{v}=a,
  \label{eq:dyn}
\end{equation}
with control bounds $|a|\le a_{\max}=0.24~\mathrm{m/s^2}$,
$|\omega|\le\omega_{\max}=0.03~\mathrm{rad/s}$, speed clipped by the environment to
$0\le v\le v_{\max}=9.5~\mathrm{m/s}$ each step (within a step $v$ may briefly reach
$v_{\max}+a_{\max}\Delta t=11.9~\mathrm{m/s}$), and decision step
$\Delta t=10~\mathrm{s}$. The ego body is the $175\times25.4$~m rectangle with
circumscribed radius $R_{\mathrm{circ}}=\tfrac12\sqrt{175^2+25.4^2}=88.4$~m and
inscribed radius $r_{\mathrm{insc}}=\tfrac12\cdot25.4=12.7$~m. Target vessels occupy
their body rectangle (or, conservatively, their circumscribed disk); unless stated
otherwise, in stand-on regimes they are extrapolated at \emph{constant velocity}
(CV). We write $p_e,p_m$ for ego/target positions,
$\bm{v}_e=v_e(\cos\theta_e,\sin\theta_e)$ for the ego velocity vector, and
$\hat{h},\hat{h}^{\perp}$ for the heading unit vector and its normal.

\subsection{Far-field regime: single-step no-collision}
\label{sec:provable-farfield}

\begin{proposition}[Far-field single-step no-collision]
\label{prop:farfield}
Let $d=\lVert p_e-p_m\rVert$ be the ego--target center distance for a single target.
If
\begin{equation}
  d \;\ge\; d_{\mathrm{safe}} + v_{\mathrm{obs,max}}\,\Delta t + \rho_{\mathrm{ego}},
  \label{eq:farfield}
\end{equation}
then after one decision step the two (circumscribed-disk) bodies cannot intersect,
i.e.\ a collision is impossible in that step. Here
$d_{\mathrm{safe}}=R_{\mathrm{circ,ego}}+R_{\mathrm{circ,obs}}$ and
$\rho_{\mathrm{ego}}=v_{\max}\Delta t+\tfrac12 a_{\max}\Delta t^2$ is the maximum
single-step ego displacement.
\end{proposition}

Substituting constants gives $R_{\mathrm{circ,ego}}=88.4$~m and, because the target
width is not in the state, a conservative square upper bound with the Table~II
safety margin $d_{\mathrm{obs,safety}}=2\ell$:
$R_{\mathrm{circ,obs}}=\tfrac12\sqrt2\,\ell+d_{\mathrm{obs,safety}}
=123.74+350=473.74$~m, so $d_{\mathrm{safe}}=562.16$~m,
$\rho_{\mathrm{ego}}=9.5\cdot10+\tfrac12\cdot0.24\cdot100=107.0$~m, and
$v_{\mathrm{obs,max}}\Delta t=95$~m, yielding a threshold of $764.16$~m
(engineering value $\sim\!780$~m with margin).

\emph{Assumptions.} (A1) single target (for $N$ targets, each $i$ must satisfy the
threshold, Proposition~\ref{prop:farfield-N}); (A2) target speed
$v_{\mathrm{obs}}\le9.5~\mathrm{m/s}$---Krasowski sets $v_{\max}$ only for the ego, so
this bound must be assumed explicitly. It is \emph{empirically verified} on the full
benchmark: across all $2000$ HandcraftedTwoVessel scenarios ($342{,}000$ obstacle
states) the maximum is $7.10~\mathrm{m/s}$, with zero states exceeding $9.5$
($\sim\!2.4~\mathrm{m/s}$ margin); higher-speed real AIS data would require raising
the threshold. (A3) $v\le v_{\max}$ at step end (environment clip); (A4) the target
is approximately CV within a step (stand-on); (A5) circumscribed-disk occupancy
(disks disjoint $\Rightarrow$ bodies disjoint).

\emph{Proof sketch.} By the triangle inequality the next-step center distance
satisfies $d_{\mathrm{next}}\ge d-\lVert\Delta p_e\rVert-\lVert\Delta p_m\rVert
\ge d-\rho_{\mathrm{ego}}-v_{\mathrm{obs,max}}\Delta t$. With \eqref{eq:farfield},
$d_{\mathrm{next}}\ge d_{\mathrm{safe}}=R_{\mathrm{circ,ego}}+R_{\mathrm{circ,obs}}$,
so the circumscribed disks---hence the bodies---are disjoint. $\square$

\begin{scope}
The triangle inequality is direction-independent, so head-on is already the true
worst case and there is no geometric gap; this is the \emph{soundest} block of the
provable layer (an elementary lemma). It is a conservative single-step
\emph{sufficient} condition---not invariance, not a near-field guarantee---and is
currently not implemented as an early exit (an $O(1)$ far-field fast path could be
added).
\end{scope}

\subsection{Give-way regime: per-step direction compliance}
\label{sec:provable-compliance}

\begin{proposition}[Per-give-way-step direction compliance]
\label{prop:compliance}
At every decision step that the state machine classifies as a give-way regime
($\rho\in\{\text{head-on},\text{crossing},\text{overtake}\}$) and for which the
projection quadratic program (QP) is feasible (\texttt{needs\_fallback}=False), the
executed safe action $u_{\mathrm{safe}}$ lies in the COLREGs-mandated compliant
half-plane (e.g.\ give-way starboard turn $\omega\le-\omega_{\mathrm{turn}}$).
\end{proposition}

\emph{Assumptions.} (B1) single target; (B2) correct state-machine classification of
$\rho$ (inheriting the $\sim\!12.6\%$ classification gap and the Lemma-1 tie-break
issue); (B3) QP feasibility (not fallback); (B4) scope $=$ give-way encounters. The
empirical give-way trigger rate is $\sim\!13$--$16\%$ \emph{at the
encounter/episode level} (measured twice over $933$ scenarios; overtake $=0$)---not
per step; stand-on violations are counted per step.

\emph{Proof.} The compliant half-plane enters the QP as a hard interval constraint
($w_{\mathrm{hi}}=-\omega_{\mathrm{turn}}$, etc.), so any projected solution
satisfies it. $\square$ (by construction; no separate theorem required).

\begin{scope}
This does \emph{not} cover the emergency regime $\rho_5$ (which occupies
$\sim\!55$--$60\%$ of conflict steps and whose \texttt{is\_emergency} predicate
structurally bypasses the compliance constraint), nor the relaxed/%
collision-minimizing fallbacks (which drop the half-plane). We therefore write only
``provably give-way-direction-compliant, per give-way step'' and \textbf{never}
``provably compliant / compliant at every step'' (the measured violation rate on the
golden 10-seed run is a nonzero $1.13$--$1.8$ per episode).
\end{scope}

\subsection{Emergency regime: a sound imminent-unavoidability certificate}
\label{sec:provable-imminent}

The emergency controller is an empirical last resort. To characterize it honestly we
do not claim it avoids collisions; instead we provide a \emph{sound sufficient
condition} that certifies when a collision is physically unavoidable, so that the
``black box'' becomes a decidable boundary. This replaces an earlier
turn-only condition that was \emph{unsound} because it ignored acceleration/%
deceleration escapes (gap\#1).

\begin{proposition}[Imminent-unavoidability certificate, single CV target]
\label{prop:imminent}
Consider a single constant-velocity (stand-on) target. Define the ego
reachable-center over-approximation $R_{\mathrm{box}}(t)$ (a heading-frame rectangle
centered at the CV extrapolation $c_e(t)=p_e+\bm{v}_e t$) with longitudinal
half-axis $B_{\mathrm{lon}}(t)=\tfrac12 a_{\mathrm{eff}}t^2$ and lateral half-axis
$L_{\mathrm{lat}}(t)$, where
\begin{align}
  a_{\mathrm{eff}} &= \sqrt{a_{\max}^2+(v_{\mathrm{bnd}}\,\omega_{\max})^2},
    \quad v_{\mathrm{bnd}}=v_{\max}+a_{\max}\Delta t=11.9~\mathrm{m/s},\\
  L_{\mathrm{lat}}(t) &=
    \begin{cases}
      \dfrac{v_{\mathrm{bnd}}}{\omega_{\max}}\bigl(1-\cos\omega_{\max}t\bigr),
        & \omega_{\max}t\le\tfrac{\pi}{2},\\[1.2ex]
      \dfrac{v_{\mathrm{bnd}}}{\omega_{\max}}
        + v_{\mathrm{bnd}}\Bigl(t-\tfrac{\pi}{2\omega_{\max}}\Bigr),
        & \text{otherwise.}
    \end{cases}
\end{align}
Let $O(t)$ be the target body rectangle (center $p_m+\bm{v}_m t$, heading
$\theta_m$). If
\begin{equation}
  \exists\, t^\ast\in[0,T]:\quad
  R_{\mathrm{box}}(t^\ast)\;\subseteq\;\bigl(O(t^\ast)\oplus \mathrm{disk}(r_{\mathrm{insc,ego}}=12.7~\mathrm{m})\bigr),
  \label{eq:imminent}
\end{equation}
then the collision is \emph{unavoidable}: every feasible control sequence collides
with the target at $t^\ast$.
\end{proposition}

\emph{Assumptions.} (A1) single target, \emph{per-pair} (per-pair unavoidable
$\Rightarrow$ global unavoidable; per-pair \emph{avoidable} does \emph{not}
compose); (A2) target CV stand-on (a maneuvering target changes $O(t)$ and voids the
guarantee---we state ``unavoidable within the CV world''); (A3) $v\le
v_{\mathrm{bnd}}=11.9$ within a step (includes the clip overshoot); (A4) the ego
inscribed disk $r_{\mathrm{insc}}=12.7~\mathrm{m}\subseteq$ body at any heading
(from the short edge; the long-edge zero-margin case is pending an independent re-check); (A5)
the certificate fires only in the region $\omega_{\max}t^\ast\le\pi/2$
($\le52.4$~s)---we verified that firing is limited by the longitudinal quadratic term
to $t_{\mathrm{crit,lon}}=\sqrt{2(\tfrac12\cdot175+12.7)/a_{\mathrm{eff}}}=21.6~\mathrm{s}<52.4~\mathrm{s}$;
(A6) single-instant containment (sufficient, not necessary); (A7) occupancy
under-approximation: $\text{obs\_width}\le$ true target width and $L_{\mathrm{obs}}\le$
true target length (a soundness precondition the caller must guarantee---passing
sizes larger than the truth over-approximates $O$ and can yield false positives).

\emph{Proof sketch (three steps).} (i) $R_{\mathrm{box}}$ over-approximates the true
reachable center: $\lVert d\bm{v}/dt\rVert=\sqrt{a^2+(v\omega)^2}\le a_{\mathrm{eff}}$,
and double integration bounds the longitudinal offset by $\tfrac12 a_{\mathrm{eff}}t^2$
while $|y'(t)|\le\int_0^t v_{\mathrm{bnd}}\sin(\omega_{\max}s)\,ds=L_{\mathrm{lat}}(t)$.
(ii) The inscribed disk under-approximates the body. (iii) If
$R_{\mathrm{box}}(t^\ast)\subseteq O(t^\ast)\oplus\mathrm{disk}(12.7)$, then every
reachable ego center lies within $12.7$~m of $O(t^\ast)$, so the ego inscribed disk
($\subseteq$ body) intersects $O(t^\ast)$: a collision. Since this holds for
\emph{all} reachable positions, every trajectory collides at $t^\ast$. $\square$

\emph{Verification.} A six-agent adversarial audit ($20{,}000$ random $+$ $735$
curated cases, with a triple-switch bang-bang escape search under exact hull-collision
geometry) found \emph{zero} true false positives; after hardening the
$v_{\mathrm{bnd}}/L_{\mathrm{lat}}$ envelope, a battery plus $1500$ fuzz cases
remained sound (0 false positives) and non-vacuous (firing at $t^\ast\le7$~s).

\begin{scope}
The certificate is a conservative \emph{one-sided sufficient} condition for
imminent unavoidability. It is \textbf{not} an early-warning cost (it is a terminal
classifier, not a planner), \textbf{not} complete or necessary (no-fire $\ne$
avoidable), and by itself certifies nothing about collision freedom. We
\textbf{never} write ``provably zero-collision'': this is an emergency-regime
classifier, not a system-wide guarantee.
\end{scope}

\subsection{Forward-invariant recursive feasibility: the provable upgrade}
\label{sec:provable-forward-invariant}

Propositions~\ref{prop:farfield}--\ref{prop:compliance} are \emph{one-step}
guarantees; a single-step no-collision constraint is not by itself recursively
feasible (a state may be safe now yet admit no safe continuation). This subsection
upgrades the shield from one-step lookahead to a \emph{forward-invariant} guarantee
on a certified clearable set, which is the paper's central provable contribution and
the differentiator against reactive CBF-QP filters (which, under a $\Delta=10$\,s
zero-order hold, carry no recursive-feasibility proof).

\paragraph{Clearable set.} Fix a single constant-velocity (CV) target. Call a finite
control sequence $m$ a \emph{certified straight-tail escape} for state $s$ if its
tail segment is \emph{constant velocity and straight} ($\omega=0\wedge a=0$) and the
sound Lipschitz clearance certificate (the same interval lower bound used for
Prop.~\ref{prop:farfield}, Section~\ref{sec:provable-farfield}) shows hull
distance $d>0$ on $[0,H]$ with a past-closest-point-of-approach (CPA) margin. Define
the \emph{clearable set} $A:=\{s:\exists\,\text{certified straight-tail escape}\}$.

\begin{lemma}[Constant-velocity tail $\Rightarrow$ permanent clearance]
\label{lem:convex-tail}
If the tail is constant velocity and straight and the target is CV, the hull distance
$g(t)=\mathrm{dist}(\text{two fixed-orientation rectangles translating at constant
relative velocity})$ is \emph{convex} in $t$; hence once $g$ is increasing past the
CPA it increases for all later $t$, so $d>0$ on $[0,H]$ together with the past-CPA
margin implies $d>0$ for all $t\ge0$.
\end{lemma}
\emph{Remark (soundness fix).} Convexity requires the tail to be constant
\emph{velocity}: $\omega=0$ alone is insufficient, since an accelerating straight tail
traces a path that is non-affine in $t$ and $g$ need not be convex. The certificate
therefore requires $a=0$ as well. [HONESTY] The certificate is a \emph{sufficient}
condition; it is sound (never certifies a colliding escape---verified by a
long-horizon fuzz with zero false positives) but conservative (it rejects some safe
escapes).

\begin{proposition}[Control invariance of the clearable set]
\label{prop:forward-invariant}
For every $s\in A$ with certified escape $m=(u_0,u_1,\dots)$, applying the first
control $u_0$ over one decision step (official semantics, with the end-of-step speed
clamp aligned to $\Delta$) yields a successor $s'=F(s,u_0)$ whose tail
$m'=(u_1,u_2,\dots)$ is a certified straight-tail escape for $s'$ (it is the second
half of the \emph{same} physical trajectory, which passed one sound certificate).
Hence $s'\in A$: $A$ is control invariant. Consequently a terminal constraint
``successor $\in A$'', realized as an $O(1)$ re-certification of one concrete backup
maneuver (the backup-maneuver terminal set $U_{\mathrm{term}}$), renders the shielded
closed loop forward invariant on $A$, i.e.\ \emph{provably collision-free for all
future time on $A$ under a single CV target}.
\end{proposition}
\emph{Proof sketch.} The tail is certified because it is a sub-arc of an already
certified trajectory (verified numerically to reproduce the original trajectory's
second half exactly under the aligned clamp). Feasibility of $U_{\mathrm{term}}$ is
non-vacuous: $u_0$ itself keeps the backup certified, so the projection QP always
retains at least the backup as a fallback. $\square$

\paragraph{Compliant forward invariance ($A\cap U_{\mathrm{colregs}}$).}
The shield must pick a give-way-\emph{compliant} first step, so the relevant invariant
is $A\cap U_{\mathrm{colregs}}$. Let the compliant backup be ``turn to the mandated
side for $t_1$\,s, then go straight.'' \textbf{Case 1} ($t_1>\Delta$): the tail still
begins with a compliant turn and remains certified by
Prop.~\ref{prop:forward-invariant}, so $s'\in A\cap U_{\mathrm{colregs}}$---this case
is proved. \textbf{Case 2} ($t_1\le\Delta$, the turn completes within one step): the
tail's first step is straight and a fresh compliant escape from $s'$ is required; a
general proof of this case is open. On the shield's real give-way state distribution,
however, Case 2 \emph{never occurs} (the slow yaw rate, $\sim\!1.7^\circ$/s for a
$175$--$260$\,m hull, forces every compliant clearing turn to exceed one decision
step): across all sampled real give-way states with a compliant backup, $100\%$ fall
in Case 1, the tail remains certified and compliant in $100\%$, and the successor
re-admits a fresh compliant certified escape in $100\%$. We therefore state:
$A\cap U_{\mathrm{colregs}}$ is forward invariant via the proved Case 1 whenever the
compliant backup's turn extends beyond one decision step, a condition met on $100\%$
of the observed give-way distribution.

\begin{scope}
[HONESTY] (i) Single CV target (multi-target is weaker, $A=\bigcap_i A_i$, and per-%
target funnels may conflict). (ii) Membership in $A$ is decided by a finite
constant-control maneuver family plus the sound-but-conservative certificate, so $A$
is an \emph{inner} approximation of the true clearable set; the residual is exactly
the $A$-membership boundary---states not clearable in \emph{any} direction---which
falls to the emergency/relaxed fallback, and is \textbf{not} a
compliance-versus-safety conflict. On the real give-way distribution this residual is
$\approx\!0$--$3\%$ (head-on/overtake $0\%$, crossing $\sim\!2.9\%$), and every
in-$A$ give-way state admits a compliant certified escape. [HONESTY] These rates were
computed with a high-accuracy local integrator that matches the official dynamics to
sub-micron; official-pipeline reconfirmation is [PENDING server, evaluation-only].
(iii) The backup-maneuver terminal constraint $U_{\mathrm{term}}$ is implemented and
unit-tested as a sound module; wiring it into the deployed shield and validating the
closed loop is [PENDING server, evaluation-only]. (iv) The give-way states are drawn
from a well-behaved shielded rollout in which holding course never collides within the
horizon, so these rates characterize the clearable-set geometry, not performance under
adversarially imminent conflict, which needs a give-way-rich adversarial source
[PENDING].
\end{scope}

\subsection{Robust extension to multiple targets}
\label{sec:provable-multi}

For $N$ targets with occupancies $O_i(t)$, collision freedom is a conjunction over
targets, $\text{no collision}\iff\bigwedge_{i=1}^{N}(\text{ego\_body}\cap O_i=\varnothing)$.
This structure is the mathematical root of the robust extension: sufficient
conditions compose as $\forall i$, while the unavoidability certificate composes as
$\exists i$.

\begin{proposition}[Far-field, multi-target]
\label{prop:farfield-N}
If $d_i\ge d_{\mathrm{safe}}+v_{\mathrm{obs,max}}\Delta t+\rho_{\mathrm{ego}}$ for
\emph{all} $i$ (the same $\sim\!780$~m threshold), then after one step the ego body
intersects no target body.
\end{proposition}
The single-target proof holds independently for each $i$; the conjunction is exactly
the definition of collision freedom, so the extension is \emph{sound and tight}
(no extra conservatism beyond the per-pair circumscribed-disk bound). The new
obligation is (A2-N): $v_{\mathrm{obs}}^i\le9.5$ for \emph{each} target---the
single-target empirical bound must be re-verified on multi-target scenarios.

\begin{proposition}[Give-way, multi-target]
\label{prop:compliance-N}
Let $G=\{i:\rho_i \text{ is give-way}\}$ with compliant half-planes $H_i$. At every
step where $(i)$ $H_G:=\bigcap_{i\in G}H_i\ne\varnothing$ (the give-way directions
coexist) and $(ii)$ the QP over $U_{\mathrm{box}}\cap H_G\cap\bigcap_i
U_{\mathrm{collision\text{-}free}}^i$ is feasible, the executed $u_{\mathrm{safe}}$ is
direction-compliant for all $i\in G$.
\end{proposition}
The half-planes enter one QP and the projected solution satisfies their
intersection. The honest scope is sharper than in the single-target case: if two
targets demand opposite give-way directions then $H_G=\varnothing$ and the step falls
to fallback (a COLREGs \emph{violation}). This is not a defect but a boundary of
COLREGs itself---a pairwise rule set with no crisp multi-vessel norm (Rule~2 /
special-circumstances); our defensible position is ``compliant when the pairwise
norms coexist, collision-avoidance-first otherwise.'' As $N$ grows the feasible set
shrinks and the give-way compliance rate \emph{decreases}; this must be measured,
not assumed.

\begin{proposition}[Imminent unavoidability, multi-target]
\label{prop:imminent-N}
With $R_{\mathrm{box}}(t)$ as in Proposition~\ref{prop:imminent} (target-independent),
if $\exists\,i,\ \exists\,t^\ast$ ($\omega_{\max}t^\ast\le\pi/2$) with
$R_{\mathrm{box}}(t^\ast)\subseteq O_i(t^\ast)\oplus\mathrm{disk}(12.7)$, then the
collision is unavoidable.
\end{proposition}
The per-pair certificate for target $i$ implies every reachable trajectory collides
with $O_i$, hence with some target; the existential $\exists i$ is exactly what the
emergency layer needs, at cost $O(N\cdot|T|)$. The half that does \emph{not} compose
is avoidability: per-pair no-fire for all $i$ does \emph{not} imply global
avoidability (evading $A$ may force the ego into the future position of $B$), so
in the multi-target world no-fire is even weaker than for a single target.

\subsection{Honest positioning of the provable layer}
\label{sec:provable-summary}

Table~\ref{tab:provable} summarizes the guarantees, their status, and their scope.
The headline claim we can defend, with all assumptions attached, is: continuous
action $\cap$ provably \emph{direction}-compliant $\cap$ COLREGs (a previously empty
intersection), plus far-field single-step provable safety ($\forall i$, sound), a
sound imminent-unavoidability certificate for the emergency regime ($\exists i$), and
---the central upgrade---\emph{forward-invariant recursive feasibility} on a
certificate-defined clearable set (Prop.~\ref{prop:forward-invariant}): the shield is
provably collision-free for all future time on that set under a single CV target, a
guarantee reactive CBF-QP filters lack under a $10$\,s hold. We do \textbf{not} claim
bare provable collision freedom, complete provable coverage, maneuvering targets, or
multi-target direction conflicts (which fall to fallback by the pairwise nature of
COLREGs); the forward-invariance is scoped to the clearable set, whose residual
($\approx\!0$--$3\%$) is deferred to fallback. The provable tier is therefore
\emph{medium-to-strong}: the invariance upgrade and gap\#1 are sound, but coverage is
still an inner approximation, the closed-loop $U_{\mathrm{term}}$ wiring and official
reconfirmation are [PENDING server, eval-only], and the empirical zero-collision
backing remains single-target until multi-vessel benchmark data are available.

\begin{table}[h]
  \caption{Provable safety layer: guarantees, status, and honest scope}
  \label{tab:provable}
  \centering
  \small
  \begin{tabular}{@{}p{0.20\linewidth}p{0.24\linewidth}p{0.46\linewidth}@{}}
    \toprule
    \textbf{Guarantee} & \textbf{Status} & \textbf{Scope / limits} \\
    \midrule
    Far-field single-step no-collision (Prop.~\ref{prop:farfield},
    \ref{prop:farfield-N})
      & Proven (elementary; single-target bound conservative, multi-target conjunction adds no extra slack)
      & $d\ge\sim\!780$~m; CV target(s); $\forall i$ conjunction; each
        $v_{\mathrm{obs}}\le9.5$; single step, not invariance; early exit not yet
        implemented. \\
    \addlinespace
    Per-give-way-step direction compliance (Prop.~\ref{prop:compliance},
    \ref{prop:compliance-N})
      & By construction (scope-limited)
      & Give-way encounters ($\sim\!13$--$16\%$, encounter-level, not per step);
        excludes emergency $\rho_5$ and fallbacks; multi-target rate drops with $N$;
        direction conflict $\to$ fallback (COLREGs pairwise boundary). Never ``every
        step.'' \\
    \addlinespace
    Imminent-unavoidability certificate (gap\#1) (Prop.~\ref{prop:imminent},
    \ref{prop:imminent-N})
      & Sound sufficient condition ($>\!20{,}000$ adversarial scenarios, 0 false positives)
      & Imminent only; CV target(s); $\exists i$ composes, avoidability does not;
        not complete / not early-warning; a classifier, not a system guarantee. \\
    \addlinespace
    Forward-invariant recursive feasibility on the clearable set $A$
    (Lem.~\ref{lem:convex-tail}, Prop.~\ref{prop:forward-invariant})
      & Proven (control invariance of $A$; certificate sound, 0 false positives on
        long-horizon fuzz). $U_{\mathrm{term}}$ module unit-tested; closed-loop wiring
        [PENDING server, eval-only]
      & Single CV target; $A$ = inner approximation (finite maneuver family +
        conservative certificate); residual $\approx\!0$--$3\%$ = states clearable in
        no direction $\to$ fallback (not a compliance conflict); $A\cap U_{\mathrm{colregs}}$
        invariance proved for Case 1 and holds on $100\%$ of the observed give-way
        distribution; numbers local-integrator, official reconfirmation [PENDING]. \\
    \bottomrule
  \end{tabular}
\end{table}

\paragraph{Emergency fallback and the failure-mechanism boundary.}
When the projection is infeasible ($P=\varnothing$) in the emergency regime
$\rho_5$, the \texttt{is\_emergency} predicate structurally bypasses the direction
constraint and Algorithm~1 minimizes collision risk---an \emph{empirical} fallback,
not a guarantee. On an adversarial purely-geometric baseline that forces the shield
into the emergency branch on $44$--$69\%$ of steps, we observe $45\%$ arrival and
two collisions; auditing confirms both occur with
\texttt{source=emergency},~$\rho=5$, i.e.\ the emergency last resort trying and
failing to avoid, and both lie outside the domain of Proposition~\ref{prop:farfield}
(far-field, non-emergency, single-step). They are not counterexamples to the headline
but they do substantiate that the emergency regime is not a guarantee. On the learned
policy (non-adversarial), we observe zero hull collisions over $10$ seeds
$\times$ $40$ episodes ($400$ total), so the ``empirically zero-collision'' statement holds; we
nonetheless \textbf{never} write provably collision-free and label the emergency
regime as an empirical fallback throughout.

%
%  诚实红线（CLAUDE.md §0 + 三份理论稿·全文必须守）：
%    绝不写 provably collision-free / provably compliant（裸）/ complete provable coverage。
%    头条能 claim = 连续动作 ∩ 可证明【方向】合规 ∩ COLREGs + 远场单步可证明 + 迫近不可避 sound 证书。
%    可证明档 = 中。到达率【不进卖点清单】（L192 B1 配平后无可证优势）。
%  TRB 硬约束：~7500 词·无附录（完整证明移交外部技术报告）。
% =====================================================================
\documentclass[12pt,letterpaper]{article}
\usepackage[utf8]{inputenc}
\usepackage{amsmath,amsthm,amssymb,bm}
\usepackage{booktabs}
\usepackage{graphicx}
\usepackage{hyperref}

\newtheorem{proposition}{Proposition}
\newtheorem{lemma}{Lemma}
\newtheorem{assumption}{Assumption}
\theoremstyle{remark}\newtheorem*{scope}{Scope and honest limits}

\title{Continuous-Action, Provably COLREGs-Direction-Compliant
Collision Avoidance for Unmanned Surface Vessels via a
Projection-Based Safety Shield}
\author{[authors]}
\date{}

\begin{document}
\maketitle

% =====================================================================
\section{Abstract}
% [DRAFT · 守诚实红线 · ~200 words target]
Provably COLREGs-\emph{direction}-compliant collision avoidance for unmanned
surface vessels (USVs) has, until now, forced a choice: either discretize the
action space and mask illegal actions (sacrificing continuous control), or keep
continuous actions but revert to soft reward shaping (sacrificing the guarantee).
We place a projection-based safety shield on maritime COLREGs directly in the
\emph{continuous} action space: a reinforcement-learning policy proposes a
continuous command, a rule state machine classifies the encounter regime and the
compliant give-way direction, and a quadratic program projects the command onto
the nearest action that is both direction-compliant and (single-step)
collision-free. We characterize exactly what the layer can and cannot prove: a
sound far-field single-step no-collision guarantee, per-give-way-step direction
compliance by construction, and a \emph{sound} imminent-unavoidability certificate
for the emergency regime. Upgrading these one-step guarantees, we prove that a
certificate-defined \emph{clearable set} is control-invariant, so a backup-maneuver
terminal constraint makes the shield \emph{forward-invariant} (recursively feasible)
on that set under a single constant-velocity target---a guarantee that reactive
CBF-QP filters lack under a $10$\,s zero-order hold---and on the shield's real
give-way state distribution a COLREGs-compliant certified escape exists at every
clearable state. We deliberately do \emph{not} claim bare provable collision freedom:
the guarantee is scoped to the clearable set (single CV), with the residual
($\approx\!0$--$3\%$, states clearable in no direction) deferred to an honestly
characterized emergency fallback. On the Krasowski--Althoff benchmark the shielded continuous policy attains
zero hull collisions over [PENDING: seeds$\times$episodes] and a structural
$\sim\!5\times$ advantage in throttle smoothness over the discrete baseline, at
comparable arrival. We report all limitations honestly, including a
training-stability collapse on a minority of seeds that we address with
policy warm-starting. [HONESTY: never "provably collision-free"; arrival is
comparable, not superior.]

% =====================================================================
\section{Introduction}
% [DRAFT skeleton]
Autonomous collision avoidance for surface vessels must satisfy the International
Regulations for Preventing Collisions at Sea (COLREGs), whose give-way rules are
\emph{directional} (e.g.\ in a head-on or crossing encounter the give-way vessel
must alter course to starboard). Two research lines have advanced maritime
collision avoidance but each concedes something essential:
\begin{itemize}
  \item \textbf{Discrete + masking} (Krasowski and Althoff 2024) obtains
  provable direction compliance by enumerating a discrete action set and masking
  non-compliant actions---at the cost of the continuous control that real actuators
  and smoothness require.
  \item \textbf{Continuous + soft reward} encodes COLREGs as a reward penalty,
  keeping continuous control but forfeiting any hard guarantee.
\end{itemize}
The intersection---continuous action $\cap$ provable direction compliance $\cap$
COLREGs---has been empty. We fill it with a projection-based safety shield adapted
from the safe-RL / action-projection literature to maritime COLREGs.

\paragraph{Contributions.}
\begin{enumerate}
  \item A projection-based COLREGs shield in the continuous action space
  (Section~\ref{sec:method}), integrated with RL in the safe-environment (SE-RL)
  style so the policy learns against the shielded dynamics.
  \item A \emph{provable safety layer} with an explicit, honestly-scoped set of
  guarantees (Section~\ref{sec:provable}): far-field single-step no-collision
  ($\forall$ target, sound), per-give-way-step direction compliance (by
  construction), and a sound imminent-unavoidability certificate for the emergency
  regime; extended to multiple targets ($\forall i$ / $\exists i$). We are explicit
  about what is \emph{not} proven.
  \item An empirical evaluation on the Krasowski--Althoff benchmark showing zero
  hull collisions and a structural smoothness advantage, with a candid account of a
  minority-seed training-stability collapse and its warm-start remedy
  (Section~\ref{sec:results}).
\end{enumerate}
% [HONESTY] We do not claim provable collision freedom or superior arrival rate.

% =====================================================================
\section{Related Work}
% [DRAFT skeleton · 待补引用]
\begin{itemize}
  \item \textbf{COLREGs-aware RL for USVs}: Krasowski and Althoff (2024) (discrete
  provably-compliant shield; our anchor and baseline); Meyer et al.\ (2020)
  (COLREG-aware RL).
  \item \textbf{Safe RL via action projection / shielding}: SE-RL and
  zonotope/QP action projection (Markgraf et al.\ 2026, serl-sprl); control
  barrier functions (CBF) and predictive safety filters as the guarantee mechanism
  we relate to.
  \item \textbf{Training-stability / exploration failure and remedies}: seed
  sensitivity in deep RL (Henderson et al.\ 2018; Agarwal et al.\ 2021, rliable);
  warm-starting / guided RL (JSRL 2022; AWAC 2020; RLPD 2023) and reverse
  curricula (Florensa et al.\ 2017), which motivate our collapse remedy.
  \item [PENDING: external baselines APF / VO-ORCA / CBF-QP for the Pareto
  comparison, Section~\ref{sec:results}.]
\end{itemize}

% =====================================================================
\section{Method}
\label{sec:method}
% [DRAFT skeleton]
\subsection{Problem setup and shielded MDP}
We adopt the Krasowski--Althoff CommonOcean setting: yaw-constrained point-mass
vessel dynamics (Eq.~\ref{eq:dyn} in Section~\ref{sec:provable}), continuous
control $u=[a,\omega]$, and the handcrafted two-vessel encounter benchmark. The
policy is trained in the SE-RL style, i.e.\ against the \emph{shielded} transition
so that it observes and adapts to the projection.

\subsection{The projection-based COLREGs shield}
At each step: (i) a rule state machine classifies the encounter regime
$\rho\in\{\text{far},\text{give-way(head-on/crossing/overtake)},\text{stand-on},
\text{emergency}\}$ and the compliant give-way half-plane; (ii) the compliant
half-plane and per-target collision-free constraints define the safe action set;
(iii) a quadratic program projects the policy command $u$ onto the nearest safe
action $u_{\mathrm{safe}}$; (iv) when the safe set is empty ($P=\varnothing$) in
the emergency regime, an empirical last-resort controller (Algorithm~1) minimizes
collision risk. [HONESTY: the emergency branch is empirical, not guaranteed.]

% ---- 可证明层正文（独立文件，含命题 1/2/3 + 多障碍 + 汇总表 + 失败机理段）----
% =====================================================================
%  TRB 2027 · Paper section draft: THE PROVABLE SAFETY LAYER
%  论文正式章节草稿（可证明层）· 由三份理论底稿忠实落稿：
%    - Paper/可证明层_正式命题_0707.md      (单障碍 命题 1/2/3)
%    - Paper/可证明层_多障碍扩展_0716.md    (多障碍 ∀i / ∃i 扩展)
%    - Paper/失败机理与紧急兜底刻画_0714.md (Limitations / 紧急兜底)
%  溯源：03 决策日志 L133-L141（原推导）· L164（诚实审计）· L165（gap#1 攻坚）·
%        L166（证书接进 repo）· L189 H（多障碍）· L192（承重订正）。
%
%  诚实红线（继承 CLAUDE.md §0 + 三份底稿）——本节任何改写都必须守：
%    绝不写 provably collision-free / provably compliant（裸）/ complete provable coverage。
%    能 claim 的只有：远场单步无碰(∀i·sound) + 每让路步方向合规(by-construction·作用域受限)
%    + 迫近不可避 sound 证书(∃i)；恒速(CV)单/多障碍模型下基准经验零船体碰撞。可证明档 = 中。
%
%  用法：% =====================================================================
%  TRB 2027 · Paper section draft: THE PROVABLE SAFETY LAYER
%  论文正式章节草稿（可证明层）· 由三份理论底稿忠实落稿：
%    - Paper/可证明层_正式命题_0707.md      (单障碍 命题 1/2/3)
%    - Paper/可证明层_多障碍扩展_0716.md    (多障碍 ∀i / ∃i 扩展)
%    - Paper/失败机理与紧急兜底刻画_0714.md (Limitations / 紧急兜底)
%  溯源：03 决策日志 L133-L141（原推导）· L164（诚实审计）· L165（gap#1 攻坚）·
%        L166（证书接进 repo）· L189 H（多障碍）· L192（承重订正）。
%
%  诚实红线（继承 CLAUDE.md §0 + 三份底稿）——本节任何改写都必须守：
%    绝不写 provably collision-free / provably compliant（裸）/ complete provable coverage。
%    能 claim 的只有：远场单步无碰(∀i·sound) + 每让路步方向合规(by-construction·作用域受限)
%    + 迫近不可避 sound 证书(∃i)；恒速(CV)单/多障碍模型下基准经验零船体碰撞。可证明档 = 中。
%
%  用法：% =====================================================================
%  TRB 2027 · Paper section draft: THE PROVABLE SAFETY LAYER
%  论文正式章节草稿（可证明层）· 由三份理论底稿忠实落稿：
%    - Paper/可证明层_正式命题_0707.md      (单障碍 命题 1/2/3)
%    - Paper/可证明层_多障碍扩展_0716.md    (多障碍 ∀i / ∃i 扩展)
%    - Paper/失败机理与紧急兜底刻画_0714.md (Limitations / 紧急兜底)
%  溯源：03 决策日志 L133-L141（原推导）· L164（诚实审计）· L165（gap#1 攻坚）·
%        L166（证书接进 repo）· L189 H（多障碍）· L192（承重订正）。
%
%  诚实红线（继承 CLAUDE.md §0 + 三份底稿）——本节任何改写都必须守：
%    绝不写 provably collision-free / provably compliant（裸）/ complete provable coverage。
%    能 claim 的只有：远场单步无碰(∀i·sound) + 每让路步方向合规(by-construction·作用域受限)
%    + 迫近不可避 sound 证书(∃i)；恒速(CV)单/多障碍模型下基准经验零船体碰撞。可证明档 = 中。
%
%  用法：\input{sec_provable_safety_layer} 进 Methods 之后（或作为独立 section）。
%  依赖宏包：\usepackage{amsmath,amsthm,amssymb,bm}。TRB 无附录/7500 词 →
%  正文留命题陈述 + 证明梗概 + 诚实作用域；完整证明移交外部技术报告（见 \cite{techreport}）。
% =====================================================================

% ---- 若主文档尚未定义，取消注释以下 theorem 环境 ----
% \newtheorem{proposition}{Proposition}
% \newtheorem{assumption}{Assumption}
% \theoremstyle{remark}\newtheorem*{scope}{Scope and honest limits}

\section{The Provable Safety Layer}
\label{sec:provable}

Our contribution is \emph{not} a claim of provable collision freedom. Prior
provably COLREGs-\emph{direction}-compliant collision avoidance either sacrificed
the continuous action space (discrete action masking, Krasowski and Althoff 2024)
or sacrificed the guarantee (continuous-action work reverts to soft reward
shaping). We place a \emph{projection-based safety shield} on maritime COLREGs in
the continuous action space and characterize \emph{exactly what can and cannot be
proven}. The safety layer is organized by encounter regime and yields three
guarantees of different strengths: a far-field single-step no-collision guarantee
(Proposition~\ref{prop:farfield}), a per-give-way-step direction-compliance
guarantee (Proposition~\ref{prop:compliance}), and a \emph{sound} imminent-%
unavoidability certificate for the emergency regime (Proposition~\ref{prop:imminent}).
Each is stated with its full assumption set, and each carries an explicit scope
boundary. Full proofs are deferred to the external technical report; here we give
statements and proof sketches. Throughout, ``proven'' means a strict lemma;
existence-level and empirical results are labeled as such.

\subsection{System setting and notation}
\label{sec:provable-setup}

We adopt the yaw-constrained point-mass vessel dynamics of Krasowski and Althoff
(2024) (vessel class SR108 container ship, $175 \times 25.4$~m):
\begin{equation}
  \dot{x}=v\cos\theta,\quad \dot{y}=v\sin\theta,\quad \dot{\theta}=\omega,\quad \dot{v}=a,
  \label{eq:dyn}
\end{equation}
with control bounds $|a|\le a_{\max}=0.24~\mathrm{m/s^2}$,
$|\omega|\le\omega_{\max}=0.03~\mathrm{rad/s}$, speed clipped by the environment to
$0\le v\le v_{\max}=9.5~\mathrm{m/s}$ each step (within a step $v$ may briefly reach
$v_{\max}+a_{\max}\Delta t=11.9~\mathrm{m/s}$), and decision step
$\Delta t=10~\mathrm{s}$. The ego body is the $175\times25.4$~m rectangle with
circumscribed radius $R_{\mathrm{circ}}=\tfrac12\sqrt{175^2+25.4^2}=88.4$~m and
inscribed radius $r_{\mathrm{insc}}=\tfrac12\cdot25.4=12.7$~m. Target vessels occupy
their body rectangle (or, conservatively, their circumscribed disk); unless stated
otherwise, in stand-on regimes they are extrapolated at \emph{constant velocity}
(CV). We write $p_e,p_m$ for ego/target positions,
$\bm{v}_e=v_e(\cos\theta_e,\sin\theta_e)$ for the ego velocity vector, and
$\hat{h},\hat{h}^{\perp}$ for the heading unit vector and its normal.

\subsection{Far-field regime: single-step no-collision}
\label{sec:provable-farfield}

\begin{proposition}[Far-field single-step no-collision]
\label{prop:farfield}
Let $d=\lVert p_e-p_m\rVert$ be the ego--target center distance for a single target.
If
\begin{equation}
  d \;\ge\; d_{\mathrm{safe}} + v_{\mathrm{obs,max}}\,\Delta t + \rho_{\mathrm{ego}},
  \label{eq:farfield}
\end{equation}
then after one decision step the two (circumscribed-disk) bodies cannot intersect,
i.e.\ a collision is impossible in that step. Here
$d_{\mathrm{safe}}=R_{\mathrm{circ,ego}}+R_{\mathrm{circ,obs}}$ and
$\rho_{\mathrm{ego}}=v_{\max}\Delta t+\tfrac12 a_{\max}\Delta t^2$ is the maximum
single-step ego displacement.
\end{proposition}

Substituting constants gives $R_{\mathrm{circ,ego}}=88.4$~m and, because the target
width is not in the state, a conservative square upper bound with the Table~II
safety margin $d_{\mathrm{obs,safety}}=2\ell$:
$R_{\mathrm{circ,obs}}=\tfrac12\sqrt2\,\ell+d_{\mathrm{obs,safety}}
=123.74+350=473.74$~m, so $d_{\mathrm{safe}}=562.16$~m,
$\rho_{\mathrm{ego}}=9.5\cdot10+\tfrac12\cdot0.24\cdot100=107.0$~m, and
$v_{\mathrm{obs,max}}\Delta t=95$~m, yielding a threshold of $764.16$~m
(engineering value $\sim\!780$~m with margin).

\emph{Assumptions.} (A1) single target (for $N$ targets, each $i$ must satisfy the
threshold, Proposition~\ref{prop:farfield-N}); (A2) target speed
$v_{\mathrm{obs}}\le9.5~\mathrm{m/s}$---Krasowski sets $v_{\max}$ only for the ego, so
this bound must be assumed explicitly. It is \emph{empirically verified} on the full
benchmark: across all $2000$ HandcraftedTwoVessel scenarios ($342{,}000$ obstacle
states) the maximum is $7.10~\mathrm{m/s}$, with zero states exceeding $9.5$
($\sim\!2.4~\mathrm{m/s}$ margin); higher-speed real AIS data would require raising
the threshold. (A3) $v\le v_{\max}$ at step end (environment clip); (A4) the target
is approximately CV within a step (stand-on); (A5) circumscribed-disk occupancy
(disks disjoint $\Rightarrow$ bodies disjoint).

\emph{Proof sketch.} By the triangle inequality the next-step center distance
satisfies $d_{\mathrm{next}}\ge d-\lVert\Delta p_e\rVert-\lVert\Delta p_m\rVert
\ge d-\rho_{\mathrm{ego}}-v_{\mathrm{obs,max}}\Delta t$. With \eqref{eq:farfield},
$d_{\mathrm{next}}\ge d_{\mathrm{safe}}=R_{\mathrm{circ,ego}}+R_{\mathrm{circ,obs}}$,
so the circumscribed disks---hence the bodies---are disjoint. $\square$

\begin{scope}
The triangle inequality is direction-independent, so head-on is already the true
worst case and there is no geometric gap; this is the \emph{soundest} block of the
provable layer (an elementary lemma). It is a conservative single-step
\emph{sufficient} condition---not invariance, not a near-field guarantee---and is
currently not implemented as an early exit (an $O(1)$ far-field fast path could be
added).
\end{scope}

\subsection{Give-way regime: per-step direction compliance}
\label{sec:provable-compliance}

\begin{proposition}[Per-give-way-step direction compliance]
\label{prop:compliance}
At every decision step that the state machine classifies as a give-way regime
($\rho\in\{\text{head-on},\text{crossing},\text{overtake}\}$) and for which the
projection quadratic program (QP) is feasible (\texttt{needs\_fallback}=False), the
executed safe action $u_{\mathrm{safe}}$ lies in the COLREGs-mandated compliant
half-plane (e.g.\ give-way starboard turn $\omega\le-\omega_{\mathrm{turn}}$).
\end{proposition}

\emph{Assumptions.} (B1) single target; (B2) correct state-machine classification of
$\rho$ (inheriting the $\sim\!12.6\%$ classification gap and the Lemma-1 tie-break
issue); (B3) QP feasibility (not fallback); (B4) scope $=$ give-way encounters. The
empirical give-way trigger rate is $\sim\!13$--$16\%$ \emph{at the
encounter/episode level} (measured twice over $933$ scenarios; overtake $=0$)---not
per step; stand-on violations are counted per step.

\emph{Proof.} The compliant half-plane enters the QP as a hard interval constraint
($w_{\mathrm{hi}}=-\omega_{\mathrm{turn}}$, etc.), so any projected solution
satisfies it. $\square$ (by construction; no separate theorem required).

\begin{scope}
This does \emph{not} cover the emergency regime $\rho_5$ (which occupies
$\sim\!55$--$60\%$ of conflict steps and whose \texttt{is\_emergency} predicate
structurally bypasses the compliance constraint), nor the relaxed/%
collision-minimizing fallbacks (which drop the half-plane). We therefore write only
``provably give-way-direction-compliant, per give-way step'' and \textbf{never}
``provably compliant / compliant at every step'' (the measured violation rate on the
golden 10-seed run is a nonzero $1.13$--$1.8$ per episode).
\end{scope}

\subsection{Emergency regime: a sound imminent-unavoidability certificate}
\label{sec:provable-imminent}

The emergency controller is an empirical last resort. To characterize it honestly we
do not claim it avoids collisions; instead we provide a \emph{sound sufficient
condition} that certifies when a collision is physically unavoidable, so that the
``black box'' becomes a decidable boundary. This replaces an earlier
turn-only condition that was \emph{unsound} because it ignored acceleration/%
deceleration escapes (gap\#1).

\begin{proposition}[Imminent-unavoidability certificate, single CV target]
\label{prop:imminent}
Consider a single constant-velocity (stand-on) target. Define the ego
reachable-center over-approximation $R_{\mathrm{box}}(t)$ (a heading-frame rectangle
centered at the CV extrapolation $c_e(t)=p_e+\bm{v}_e t$) with longitudinal
half-axis $B_{\mathrm{lon}}(t)=\tfrac12 a_{\mathrm{eff}}t^2$ and lateral half-axis
$L_{\mathrm{lat}}(t)$, where
\begin{align}
  a_{\mathrm{eff}} &= \sqrt{a_{\max}^2+(v_{\mathrm{bnd}}\,\omega_{\max})^2},
    \quad v_{\mathrm{bnd}}=v_{\max}+a_{\max}\Delta t=11.9~\mathrm{m/s},\\
  L_{\mathrm{lat}}(t) &=
    \begin{cases}
      \dfrac{v_{\mathrm{bnd}}}{\omega_{\max}}\bigl(1-\cos\omega_{\max}t\bigr),
        & \omega_{\max}t\le\tfrac{\pi}{2},\\[1.2ex]
      \dfrac{v_{\mathrm{bnd}}}{\omega_{\max}}
        + v_{\mathrm{bnd}}\Bigl(t-\tfrac{\pi}{2\omega_{\max}}\Bigr),
        & \text{otherwise.}
    \end{cases}
\end{align}
Let $O(t)$ be the target body rectangle (center $p_m+\bm{v}_m t$, heading
$\theta_m$). If
\begin{equation}
  \exists\, t^\ast\in[0,T]:\quad
  R_{\mathrm{box}}(t^\ast)\;\subseteq\;\bigl(O(t^\ast)\oplus \mathrm{disk}(r_{\mathrm{insc,ego}}=12.7~\mathrm{m})\bigr),
  \label{eq:imminent}
\end{equation}
then the collision is \emph{unavoidable}: every feasible control sequence collides
with the target at $t^\ast$.
\end{proposition}

\emph{Assumptions.} (A1) single target, \emph{per-pair} (per-pair unavoidable
$\Rightarrow$ global unavoidable; per-pair \emph{avoidable} does \emph{not}
compose); (A2) target CV stand-on (a maneuvering target changes $O(t)$ and voids the
guarantee---we state ``unavoidable within the CV world''); (A3) $v\le
v_{\mathrm{bnd}}=11.9$ within a step (includes the clip overshoot); (A4) the ego
inscribed disk $r_{\mathrm{insc}}=12.7~\mathrm{m}\subseteq$ body at any heading
(from the short edge; the long-edge zero-margin case is pending an independent re-check); (A5)
the certificate fires only in the region $\omega_{\max}t^\ast\le\pi/2$
($\le52.4$~s)---we verified that firing is limited by the longitudinal quadratic term
to $t_{\mathrm{crit,lon}}=\sqrt{2(\tfrac12\cdot175+12.7)/a_{\mathrm{eff}}}=21.6~\mathrm{s}<52.4~\mathrm{s}$;
(A6) single-instant containment (sufficient, not necessary); (A7) occupancy
under-approximation: $\text{obs\_width}\le$ true target width and $L_{\mathrm{obs}}\le$
true target length (a soundness precondition the caller must guarantee---passing
sizes larger than the truth over-approximates $O$ and can yield false positives).

\emph{Proof sketch (three steps).} (i) $R_{\mathrm{box}}$ over-approximates the true
reachable center: $\lVert d\bm{v}/dt\rVert=\sqrt{a^2+(v\omega)^2}\le a_{\mathrm{eff}}$,
and double integration bounds the longitudinal offset by $\tfrac12 a_{\mathrm{eff}}t^2$
while $|y'(t)|\le\int_0^t v_{\mathrm{bnd}}\sin(\omega_{\max}s)\,ds=L_{\mathrm{lat}}(t)$.
(ii) The inscribed disk under-approximates the body. (iii) If
$R_{\mathrm{box}}(t^\ast)\subseteq O(t^\ast)\oplus\mathrm{disk}(12.7)$, then every
reachable ego center lies within $12.7$~m of $O(t^\ast)$, so the ego inscribed disk
($\subseteq$ body) intersects $O(t^\ast)$: a collision. Since this holds for
\emph{all} reachable positions, every trajectory collides at $t^\ast$. $\square$

\emph{Verification.} A six-agent adversarial audit ($20{,}000$ random $+$ $735$
curated cases, with a triple-switch bang-bang escape search under exact hull-collision
geometry) found \emph{zero} true false positives; after hardening the
$v_{\mathrm{bnd}}/L_{\mathrm{lat}}$ envelope, a battery plus $1500$ fuzz cases
remained sound (0 false positives) and non-vacuous (firing at $t^\ast\le7$~s).

\begin{scope}
The certificate is a conservative \emph{one-sided sufficient} condition for
imminent unavoidability. It is \textbf{not} an early-warning cost (it is a terminal
classifier, not a planner), \textbf{not} complete or necessary (no-fire $\ne$
avoidable), and by itself certifies nothing about collision freedom. We
\textbf{never} write ``provably zero-collision'': this is an emergency-regime
classifier, not a system-wide guarantee.
\end{scope}

\subsection{Forward-invariant recursive feasibility: the provable upgrade}
\label{sec:provable-forward-invariant}

Propositions~\ref{prop:farfield}--\ref{prop:compliance} are \emph{one-step}
guarantees; a single-step no-collision constraint is not by itself recursively
feasible (a state may be safe now yet admit no safe continuation). This subsection
upgrades the shield from one-step lookahead to a \emph{forward-invariant} guarantee
on a certified clearable set, which is the paper's central provable contribution and
the differentiator against reactive CBF-QP filters (which, under a $\Delta=10$\,s
zero-order hold, carry no recursive-feasibility proof).

\paragraph{Clearable set.} Fix a single constant-velocity (CV) target. Call a finite
control sequence $m$ a \emph{certified straight-tail escape} for state $s$ if its
tail segment is \emph{constant velocity and straight} ($\omega=0\wedge a=0$) and the
sound Lipschitz clearance certificate (the same interval lower bound used for
Prop.~\ref{prop:farfield}, Section~\ref{sec:provable-farfield}) shows hull
distance $d>0$ on $[0,H]$ with a past-closest-point-of-approach (CPA) margin. Define
the \emph{clearable set} $A:=\{s:\exists\,\text{certified straight-tail escape}\}$.

\begin{lemma}[Constant-velocity tail $\Rightarrow$ permanent clearance]
\label{lem:convex-tail}
If the tail is constant velocity and straight and the target is CV, the hull distance
$g(t)=\mathrm{dist}(\text{two fixed-orientation rectangles translating at constant
relative velocity})$ is \emph{convex} in $t$; hence once $g$ is increasing past the
CPA it increases for all later $t$, so $d>0$ on $[0,H]$ together with the past-CPA
margin implies $d>0$ for all $t\ge0$.
\end{lemma}
\emph{Remark (soundness fix).} Convexity requires the tail to be constant
\emph{velocity}: $\omega=0$ alone is insufficient, since an accelerating straight tail
traces a path that is non-affine in $t$ and $g$ need not be convex. The certificate
therefore requires $a=0$ as well. [HONESTY] The certificate is a \emph{sufficient}
condition; it is sound (never certifies a colliding escape---verified by a
long-horizon fuzz with zero false positives) but conservative (it rejects some safe
escapes).

\begin{proposition}[Control invariance of the clearable set]
\label{prop:forward-invariant}
For every $s\in A$ with certified escape $m=(u_0,u_1,\dots)$, applying the first
control $u_0$ over one decision step (official semantics, with the end-of-step speed
clamp aligned to $\Delta$) yields a successor $s'=F(s,u_0)$ whose tail
$m'=(u_1,u_2,\dots)$ is a certified straight-tail escape for $s'$ (it is the second
half of the \emph{same} physical trajectory, which passed one sound certificate).
Hence $s'\in A$: $A$ is control invariant. Consequently a terminal constraint
``successor $\in A$'', realized as an $O(1)$ re-certification of one concrete backup
maneuver (the backup-maneuver terminal set $U_{\mathrm{term}}$), renders the shielded
closed loop forward invariant on $A$, i.e.\ \emph{provably collision-free for all
future time on $A$ under a single CV target}.
\end{proposition}
\emph{Proof sketch.} The tail is certified because it is a sub-arc of an already
certified trajectory (verified numerically to reproduce the original trajectory's
second half exactly under the aligned clamp). Feasibility of $U_{\mathrm{term}}$ is
non-vacuous: $u_0$ itself keeps the backup certified, so the projection QP always
retains at least the backup as a fallback. $\square$

\paragraph{Compliant forward invariance ($A\cap U_{\mathrm{colregs}}$).}
The shield must pick a give-way-\emph{compliant} first step, so the relevant invariant
is $A\cap U_{\mathrm{colregs}}$. Let the compliant backup be ``turn to the mandated
side for $t_1$\,s, then go straight.'' \textbf{Case 1} ($t_1>\Delta$): the tail still
begins with a compliant turn and remains certified by
Prop.~\ref{prop:forward-invariant}, so $s'\in A\cap U_{\mathrm{colregs}}$---this case
is proved. \textbf{Case 2} ($t_1\le\Delta$, the turn completes within one step): the
tail's first step is straight and a fresh compliant escape from $s'$ is required; a
general proof of this case is open. On the shield's real give-way state distribution,
however, Case 2 \emph{never occurs} (the slow yaw rate, $\sim\!1.7^\circ$/s for a
$175$--$260$\,m hull, forces every compliant clearing turn to exceed one decision
step): across all sampled real give-way states with a compliant backup, $100\%$ fall
in Case 1, the tail remains certified and compliant in $100\%$, and the successor
re-admits a fresh compliant certified escape in $100\%$. We therefore state:
$A\cap U_{\mathrm{colregs}}$ is forward invariant via the proved Case 1 whenever the
compliant backup's turn extends beyond one decision step, a condition met on $100\%$
of the observed give-way distribution.

\begin{scope}
[HONESTY] (i) Single CV target (multi-target is weaker, $A=\bigcap_i A_i$, and per-%
target funnels may conflict). (ii) Membership in $A$ is decided by a finite
constant-control maneuver family plus the sound-but-conservative certificate, so $A$
is an \emph{inner} approximation of the true clearable set; the residual is exactly
the $A$-membership boundary---states not clearable in \emph{any} direction---which
falls to the emergency/relaxed fallback, and is \textbf{not} a
compliance-versus-safety conflict. On the real give-way distribution this residual is
$\approx\!0$--$3\%$ (head-on/overtake $0\%$, crossing $\sim\!2.9\%$), and every
in-$A$ give-way state admits a compliant certified escape. [HONESTY] These rates were
computed with a high-accuracy local integrator that matches the official dynamics to
sub-micron; official-pipeline reconfirmation is [PENDING server, evaluation-only].
(iii) The backup-maneuver terminal constraint $U_{\mathrm{term}}$ is implemented and
unit-tested as a sound module; wiring it into the deployed shield and validating the
closed loop is [PENDING server, evaluation-only]. (iv) The give-way states are drawn
from a well-behaved shielded rollout in which holding course never collides within the
horizon, so these rates characterize the clearable-set geometry, not performance under
adversarially imminent conflict, which needs a give-way-rich adversarial source
[PENDING].
\end{scope}

\subsection{Robust extension to multiple targets}
\label{sec:provable-multi}

For $N$ targets with occupancies $O_i(t)$, collision freedom is a conjunction over
targets, $\text{no collision}\iff\bigwedge_{i=1}^{N}(\text{ego\_body}\cap O_i=\varnothing)$.
This structure is the mathematical root of the robust extension: sufficient
conditions compose as $\forall i$, while the unavoidability certificate composes as
$\exists i$.

\begin{proposition}[Far-field, multi-target]
\label{prop:farfield-N}
If $d_i\ge d_{\mathrm{safe}}+v_{\mathrm{obs,max}}\Delta t+\rho_{\mathrm{ego}}$ for
\emph{all} $i$ (the same $\sim\!780$~m threshold), then after one step the ego body
intersects no target body.
\end{proposition}
The single-target proof holds independently for each $i$; the conjunction is exactly
the definition of collision freedom, so the extension is \emph{sound and tight}
(no extra conservatism beyond the per-pair circumscribed-disk bound). The new
obligation is (A2-N): $v_{\mathrm{obs}}^i\le9.5$ for \emph{each} target---the
single-target empirical bound must be re-verified on multi-target scenarios.

\begin{proposition}[Give-way, multi-target]
\label{prop:compliance-N}
Let $G=\{i:\rho_i \text{ is give-way}\}$ with compliant half-planes $H_i$. At every
step where $(i)$ $H_G:=\bigcap_{i\in G}H_i\ne\varnothing$ (the give-way directions
coexist) and $(ii)$ the QP over $U_{\mathrm{box}}\cap H_G\cap\bigcap_i
U_{\mathrm{collision\text{-}free}}^i$ is feasible, the executed $u_{\mathrm{safe}}$ is
direction-compliant for all $i\in G$.
\end{proposition}
The half-planes enter one QP and the projected solution satisfies their
intersection. The honest scope is sharper than in the single-target case: if two
targets demand opposite give-way directions then $H_G=\varnothing$ and the step falls
to fallback (a COLREGs \emph{violation}). This is not a defect but a boundary of
COLREGs itself---a pairwise rule set with no crisp multi-vessel norm (Rule~2 /
special-circumstances); our defensible position is ``compliant when the pairwise
norms coexist, collision-avoidance-first otherwise.'' As $N$ grows the feasible set
shrinks and the give-way compliance rate \emph{decreases}; this must be measured,
not assumed.

\begin{proposition}[Imminent unavoidability, multi-target]
\label{prop:imminent-N}
With $R_{\mathrm{box}}(t)$ as in Proposition~\ref{prop:imminent} (target-independent),
if $\exists\,i,\ \exists\,t^\ast$ ($\omega_{\max}t^\ast\le\pi/2$) with
$R_{\mathrm{box}}(t^\ast)\subseteq O_i(t^\ast)\oplus\mathrm{disk}(12.7)$, then the
collision is unavoidable.
\end{proposition}
The per-pair certificate for target $i$ implies every reachable trajectory collides
with $O_i$, hence with some target; the existential $\exists i$ is exactly what the
emergency layer needs, at cost $O(N\cdot|T|)$. The half that does \emph{not} compose
is avoidability: per-pair no-fire for all $i$ does \emph{not} imply global
avoidability (evading $A$ may force the ego into the future position of $B$), so
in the multi-target world no-fire is even weaker than for a single target.

\subsection{Honest positioning of the provable layer}
\label{sec:provable-summary}

Table~\ref{tab:provable} summarizes the guarantees, their status, and their scope.
The headline claim we can defend, with all assumptions attached, is: continuous
action $\cap$ provably \emph{direction}-compliant $\cap$ COLREGs (a previously empty
intersection), plus far-field single-step provable safety ($\forall i$, sound), a
sound imminent-unavoidability certificate for the emergency regime ($\exists i$), and
---the central upgrade---\emph{forward-invariant recursive feasibility} on a
certificate-defined clearable set (Prop.~\ref{prop:forward-invariant}): the shield is
provably collision-free for all future time on that set under a single CV target, a
guarantee reactive CBF-QP filters lack under a $10$\,s hold. We do \textbf{not} claim
bare provable collision freedom, complete provable coverage, maneuvering targets, or
multi-target direction conflicts (which fall to fallback by the pairwise nature of
COLREGs); the forward-invariance is scoped to the clearable set, whose residual
($\approx\!0$--$3\%$) is deferred to fallback. The provable tier is therefore
\emph{medium-to-strong}: the invariance upgrade and gap\#1 are sound, but coverage is
still an inner approximation, the closed-loop $U_{\mathrm{term}}$ wiring and official
reconfirmation are [PENDING server, eval-only], and the empirical zero-collision
backing remains single-target until multi-vessel benchmark data are available.

\begin{table}[h]
  \caption{Provable safety layer: guarantees, status, and honest scope}
  \label{tab:provable}
  \centering
  \small
  \begin{tabular}{@{}p{0.20\linewidth}p{0.24\linewidth}p{0.46\linewidth}@{}}
    \toprule
    \textbf{Guarantee} & \textbf{Status} & \textbf{Scope / limits} \\
    \midrule
    Far-field single-step no-collision (Prop.~\ref{prop:farfield},
    \ref{prop:farfield-N})
      & Proven (elementary; single-target bound conservative, multi-target conjunction adds no extra slack)
      & $d\ge\sim\!780$~m; CV target(s); $\forall i$ conjunction; each
        $v_{\mathrm{obs}}\le9.5$; single step, not invariance; early exit not yet
        implemented. \\
    \addlinespace
    Per-give-way-step direction compliance (Prop.~\ref{prop:compliance},
    \ref{prop:compliance-N})
      & By construction (scope-limited)
      & Give-way encounters ($\sim\!13$--$16\%$, encounter-level, not per step);
        excludes emergency $\rho_5$ and fallbacks; multi-target rate drops with $N$;
        direction conflict $\to$ fallback (COLREGs pairwise boundary). Never ``every
        step.'' \\
    \addlinespace
    Imminent-unavoidability certificate (gap\#1) (Prop.~\ref{prop:imminent},
    \ref{prop:imminent-N})
      & Sound sufficient condition ($>\!20{,}000$ adversarial scenarios, 0 false positives)
      & Imminent only; CV target(s); $\exists i$ composes, avoidability does not;
        not complete / not early-warning; a classifier, not a system guarantee. \\
    \addlinespace
    Forward-invariant recursive feasibility on the clearable set $A$
    (Lem.~\ref{lem:convex-tail}, Prop.~\ref{prop:forward-invariant})
      & Proven (control invariance of $A$; certificate sound, 0 false positives on
        long-horizon fuzz). $U_{\mathrm{term}}$ module unit-tested; closed-loop wiring
        [PENDING server, eval-only]
      & Single CV target; $A$ = inner approximation (finite maneuver family +
        conservative certificate); residual $\approx\!0$--$3\%$ = states clearable in
        no direction $\to$ fallback (not a compliance conflict); $A\cap U_{\mathrm{colregs}}$
        invariance proved for Case 1 and holds on $100\%$ of the observed give-way
        distribution; numbers local-integrator, official reconfirmation [PENDING]. \\
    \bottomrule
  \end{tabular}
\end{table}

\paragraph{Emergency fallback and the failure-mechanism boundary.}
When the projection is infeasible ($P=\varnothing$) in the emergency regime
$\rho_5$, the \texttt{is\_emergency} predicate structurally bypasses the direction
constraint and Algorithm~1 minimizes collision risk---an \emph{empirical} fallback,
not a guarantee. On an adversarial purely-geometric baseline that forces the shield
into the emergency branch on $44$--$69\%$ of steps, we observe $45\%$ arrival and
two collisions; auditing confirms both occur with
\texttt{source=emergency},~$\rho=5$, i.e.\ the emergency last resort trying and
failing to avoid, and both lie outside the domain of Proposition~\ref{prop:farfield}
(far-field, non-emergency, single-step). They are not counterexamples to the headline
but they do substantiate that the emergency regime is not a guarantee. On the learned
policy (non-adversarial), we observe zero hull collisions over $10$ seeds
$\times$ $40$ episodes ($400$ total), so the ``empirically zero-collision'' statement holds; we
nonetheless \textbf{never} write provably collision-free and label the emergency
regime as an empirical fallback throughout.
 进 Methods 之后（或作为独立 section）。
%  依赖宏包：\usepackage{amsmath,amsthm,amssymb,bm}。TRB 无附录/7500 词 →
%  正文留命题陈述 + 证明梗概 + 诚实作用域；完整证明移交外部技术报告（见 \cite{techreport}）。
% =====================================================================

% ---- 若主文档尚未定义，取消注释以下 theorem 环境 ----
% \newtheorem{proposition}{Proposition}
% \newtheorem{assumption}{Assumption}
% \theoremstyle{remark}\newtheorem*{scope}{Scope and honest limits}

\section{The Provable Safety Layer}
\label{sec:provable}

Our contribution is \emph{not} a claim of provable collision freedom. Prior
provably COLREGs-\emph{direction}-compliant collision avoidance either sacrificed
the continuous action space (discrete action masking, Krasowski and Althoff 2024)
or sacrificed the guarantee (continuous-action work reverts to soft reward
shaping). We place a \emph{projection-based safety shield} on maritime COLREGs in
the continuous action space and characterize \emph{exactly what can and cannot be
proven}. The safety layer is organized by encounter regime and yields three
guarantees of different strengths: a far-field single-step no-collision guarantee
(Proposition~\ref{prop:farfield}), a per-give-way-step direction-compliance
guarantee (Proposition~\ref{prop:compliance}), and a \emph{sound} imminent-%
unavoidability certificate for the emergency regime (Proposition~\ref{prop:imminent}).
Each is stated with its full assumption set, and each carries an explicit scope
boundary. Full proofs are deferred to the external technical report; here we give
statements and proof sketches. Throughout, ``proven'' means a strict lemma;
existence-level and empirical results are labeled as such.

\subsection{System setting and notation}
\label{sec:provable-setup}

We adopt the yaw-constrained point-mass vessel dynamics of Krasowski and Althoff
(2024) (vessel class SR108 container ship, $175 \times 25.4$~m):
\begin{equation}
  \dot{x}=v\cos\theta,\quad \dot{y}=v\sin\theta,\quad \dot{\theta}=\omega,\quad \dot{v}=a,
  \label{eq:dyn}
\end{equation}
with control bounds $|a|\le a_{\max}=0.24~\mathrm{m/s^2}$,
$|\omega|\le\omega_{\max}=0.03~\mathrm{rad/s}$, speed clipped by the environment to
$0\le v\le v_{\max}=9.5~\mathrm{m/s}$ each step (within a step $v$ may briefly reach
$v_{\max}+a_{\max}\Delta t=11.9~\mathrm{m/s}$), and decision step
$\Delta t=10~\mathrm{s}$. The ego body is the $175\times25.4$~m rectangle with
circumscribed radius $R_{\mathrm{circ}}=\tfrac12\sqrt{175^2+25.4^2}=88.4$~m and
inscribed radius $r_{\mathrm{insc}}=\tfrac12\cdot25.4=12.7$~m. Target vessels occupy
their body rectangle (or, conservatively, their circumscribed disk); unless stated
otherwise, in stand-on regimes they are extrapolated at \emph{constant velocity}
(CV). We write $p_e,p_m$ for ego/target positions,
$\bm{v}_e=v_e(\cos\theta_e,\sin\theta_e)$ for the ego velocity vector, and
$\hat{h},\hat{h}^{\perp}$ for the heading unit vector and its normal.

\subsection{Far-field regime: single-step no-collision}
\label{sec:provable-farfield}

\begin{proposition}[Far-field single-step no-collision]
\label{prop:farfield}
Let $d=\lVert p_e-p_m\rVert$ be the ego--target center distance for a single target.
If
\begin{equation}
  d \;\ge\; d_{\mathrm{safe}} + v_{\mathrm{obs,max}}\,\Delta t + \rho_{\mathrm{ego}},
  \label{eq:farfield}
\end{equation}
then after one decision step the two (circumscribed-disk) bodies cannot intersect,
i.e.\ a collision is impossible in that step. Here
$d_{\mathrm{safe}}=R_{\mathrm{circ,ego}}+R_{\mathrm{circ,obs}}$ and
$\rho_{\mathrm{ego}}=v_{\max}\Delta t+\tfrac12 a_{\max}\Delta t^2$ is the maximum
single-step ego displacement.
\end{proposition}

Substituting constants gives $R_{\mathrm{circ,ego}}=88.4$~m and, because the target
width is not in the state, a conservative square upper bound with the Table~II
safety margin $d_{\mathrm{obs,safety}}=2\ell$:
$R_{\mathrm{circ,obs}}=\tfrac12\sqrt2\,\ell+d_{\mathrm{obs,safety}}
=123.74+350=473.74$~m, so $d_{\mathrm{safe}}=562.16$~m,
$\rho_{\mathrm{ego}}=9.5\cdot10+\tfrac12\cdot0.24\cdot100=107.0$~m, and
$v_{\mathrm{obs,max}}\Delta t=95$~m, yielding a threshold of $764.16$~m
(engineering value $\sim\!780$~m with margin).

\emph{Assumptions.} (A1) single target (for $N$ targets, each $i$ must satisfy the
threshold, Proposition~\ref{prop:farfield-N}); (A2) target speed
$v_{\mathrm{obs}}\le9.5~\mathrm{m/s}$---Krasowski sets $v_{\max}$ only for the ego, so
this bound must be assumed explicitly. It is \emph{empirically verified} on the full
benchmark: across all $2000$ HandcraftedTwoVessel scenarios ($342{,}000$ obstacle
states) the maximum is $7.10~\mathrm{m/s}$, with zero states exceeding $9.5$
($\sim\!2.4~\mathrm{m/s}$ margin); higher-speed real AIS data would require raising
the threshold. (A3) $v\le v_{\max}$ at step end (environment clip); (A4) the target
is approximately CV within a step (stand-on); (A5) circumscribed-disk occupancy
(disks disjoint $\Rightarrow$ bodies disjoint).

\emph{Proof sketch.} By the triangle inequality the next-step center distance
satisfies $d_{\mathrm{next}}\ge d-\lVert\Delta p_e\rVert-\lVert\Delta p_m\rVert
\ge d-\rho_{\mathrm{ego}}-v_{\mathrm{obs,max}}\Delta t$. With \eqref{eq:farfield},
$d_{\mathrm{next}}\ge d_{\mathrm{safe}}=R_{\mathrm{circ,ego}}+R_{\mathrm{circ,obs}}$,
so the circumscribed disks---hence the bodies---are disjoint. $\square$

\begin{scope}
The triangle inequality is direction-independent, so head-on is already the true
worst case and there is no geometric gap; this is the \emph{soundest} block of the
provable layer (an elementary lemma). It is a conservative single-step
\emph{sufficient} condition---not invariance, not a near-field guarantee---and is
currently not implemented as an early exit (an $O(1)$ far-field fast path could be
added).
\end{scope}

\subsection{Give-way regime: per-step direction compliance}
\label{sec:provable-compliance}

\begin{proposition}[Per-give-way-step direction compliance]
\label{prop:compliance}
At every decision step that the state machine classifies as a give-way regime
($\rho\in\{\text{head-on},\text{crossing},\text{overtake}\}$) and for which the
projection quadratic program (QP) is feasible (\texttt{needs\_fallback}=False), the
executed safe action $u_{\mathrm{safe}}$ lies in the COLREGs-mandated compliant
half-plane (e.g.\ give-way starboard turn $\omega\le-\omega_{\mathrm{turn}}$).
\end{proposition}

\emph{Assumptions.} (B1) single target; (B2) correct state-machine classification of
$\rho$ (inheriting the $\sim\!12.6\%$ classification gap and the Lemma-1 tie-break
issue); (B3) QP feasibility (not fallback); (B4) scope $=$ give-way encounters. The
empirical give-way trigger rate is $\sim\!13$--$16\%$ \emph{at the
encounter/episode level} (measured twice over $933$ scenarios; overtake $=0$)---not
per step; stand-on violations are counted per step.

\emph{Proof.} The compliant half-plane enters the QP as a hard interval constraint
($w_{\mathrm{hi}}=-\omega_{\mathrm{turn}}$, etc.), so any projected solution
satisfies it. $\square$ (by construction; no separate theorem required).

\begin{scope}
This does \emph{not} cover the emergency regime $\rho_5$ (which occupies
$\sim\!55$--$60\%$ of conflict steps and whose \texttt{is\_emergency} predicate
structurally bypasses the compliance constraint), nor the relaxed/%
collision-minimizing fallbacks (which drop the half-plane). We therefore write only
``provably give-way-direction-compliant, per give-way step'' and \textbf{never}
``provably compliant / compliant at every step'' (the measured violation rate on the
golden 10-seed run is a nonzero $1.13$--$1.8$ per episode).
\end{scope}

\subsection{Emergency regime: a sound imminent-unavoidability certificate}
\label{sec:provable-imminent}

The emergency controller is an empirical last resort. To characterize it honestly we
do not claim it avoids collisions; instead we provide a \emph{sound sufficient
condition} that certifies when a collision is physically unavoidable, so that the
``black box'' becomes a decidable boundary. This replaces an earlier
turn-only condition that was \emph{unsound} because it ignored acceleration/%
deceleration escapes (gap\#1).

\begin{proposition}[Imminent-unavoidability certificate, single CV target]
\label{prop:imminent}
Consider a single constant-velocity (stand-on) target. Define the ego
reachable-center over-approximation $R_{\mathrm{box}}(t)$ (a heading-frame rectangle
centered at the CV extrapolation $c_e(t)=p_e+\bm{v}_e t$) with longitudinal
half-axis $B_{\mathrm{lon}}(t)=\tfrac12 a_{\mathrm{eff}}t^2$ and lateral half-axis
$L_{\mathrm{lat}}(t)$, where
\begin{align}
  a_{\mathrm{eff}} &= \sqrt{a_{\max}^2+(v_{\mathrm{bnd}}\,\omega_{\max})^2},
    \quad v_{\mathrm{bnd}}=v_{\max}+a_{\max}\Delta t=11.9~\mathrm{m/s},\\
  L_{\mathrm{lat}}(t) &=
    \begin{cases}
      \dfrac{v_{\mathrm{bnd}}}{\omega_{\max}}\bigl(1-\cos\omega_{\max}t\bigr),
        & \omega_{\max}t\le\tfrac{\pi}{2},\\[1.2ex]
      \dfrac{v_{\mathrm{bnd}}}{\omega_{\max}}
        + v_{\mathrm{bnd}}\Bigl(t-\tfrac{\pi}{2\omega_{\max}}\Bigr),
        & \text{otherwise.}
    \end{cases}
\end{align}
Let $O(t)$ be the target body rectangle (center $p_m+\bm{v}_m t$, heading
$\theta_m$). If
\begin{equation}
  \exists\, t^\ast\in[0,T]:\quad
  R_{\mathrm{box}}(t^\ast)\;\subseteq\;\bigl(O(t^\ast)\oplus \mathrm{disk}(r_{\mathrm{insc,ego}}=12.7~\mathrm{m})\bigr),
  \label{eq:imminent}
\end{equation}
then the collision is \emph{unavoidable}: every feasible control sequence collides
with the target at $t^\ast$.
\end{proposition}

\emph{Assumptions.} (A1) single target, \emph{per-pair} (per-pair unavoidable
$\Rightarrow$ global unavoidable; per-pair \emph{avoidable} does \emph{not}
compose); (A2) target CV stand-on (a maneuvering target changes $O(t)$ and voids the
guarantee---we state ``unavoidable within the CV world''); (A3) $v\le
v_{\mathrm{bnd}}=11.9$ within a step (includes the clip overshoot); (A4) the ego
inscribed disk $r_{\mathrm{insc}}=12.7~\mathrm{m}\subseteq$ body at any heading
(from the short edge; the long-edge zero-margin case is pending an independent re-check); (A5)
the certificate fires only in the region $\omega_{\max}t^\ast\le\pi/2$
($\le52.4$~s)---we verified that firing is limited by the longitudinal quadratic term
to $t_{\mathrm{crit,lon}}=\sqrt{2(\tfrac12\cdot175+12.7)/a_{\mathrm{eff}}}=21.6~\mathrm{s}<52.4~\mathrm{s}$;
(A6) single-instant containment (sufficient, not necessary); (A7) occupancy
under-approximation: $\text{obs\_width}\le$ true target width and $L_{\mathrm{obs}}\le$
true target length (a soundness precondition the caller must guarantee---passing
sizes larger than the truth over-approximates $O$ and can yield false positives).

\emph{Proof sketch (three steps).} (i) $R_{\mathrm{box}}$ over-approximates the true
reachable center: $\lVert d\bm{v}/dt\rVert=\sqrt{a^2+(v\omega)^2}\le a_{\mathrm{eff}}$,
and double integration bounds the longitudinal offset by $\tfrac12 a_{\mathrm{eff}}t^2$
while $|y'(t)|\le\int_0^t v_{\mathrm{bnd}}\sin(\omega_{\max}s)\,ds=L_{\mathrm{lat}}(t)$.
(ii) The inscribed disk under-approximates the body. (iii) If
$R_{\mathrm{box}}(t^\ast)\subseteq O(t^\ast)\oplus\mathrm{disk}(12.7)$, then every
reachable ego center lies within $12.7$~m of $O(t^\ast)$, so the ego inscribed disk
($\subseteq$ body) intersects $O(t^\ast)$: a collision. Since this holds for
\emph{all} reachable positions, every trajectory collides at $t^\ast$. $\square$

\emph{Verification.} A six-agent adversarial audit ($20{,}000$ random $+$ $735$
curated cases, with a triple-switch bang-bang escape search under exact hull-collision
geometry) found \emph{zero} true false positives; after hardening the
$v_{\mathrm{bnd}}/L_{\mathrm{lat}}$ envelope, a battery plus $1500$ fuzz cases
remained sound (0 false positives) and non-vacuous (firing at $t^\ast\le7$~s).

\begin{scope}
The certificate is a conservative \emph{one-sided sufficient} condition for
imminent unavoidability. It is \textbf{not} an early-warning cost (it is a terminal
classifier, not a planner), \textbf{not} complete or necessary (no-fire $\ne$
avoidable), and by itself certifies nothing about collision freedom. We
\textbf{never} write ``provably zero-collision'': this is an emergency-regime
classifier, not a system-wide guarantee.
\end{scope}

\subsection{Forward-invariant recursive feasibility: the provable upgrade}
\label{sec:provable-forward-invariant}

Propositions~\ref{prop:farfield}--\ref{prop:compliance} are \emph{one-step}
guarantees; a single-step no-collision constraint is not by itself recursively
feasible (a state may be safe now yet admit no safe continuation). This subsection
upgrades the shield from one-step lookahead to a \emph{forward-invariant} guarantee
on a certified clearable set, which is the paper's central provable contribution and
the differentiator against reactive CBF-QP filters (which, under a $\Delta=10$\,s
zero-order hold, carry no recursive-feasibility proof).

\paragraph{Clearable set.} Fix a single constant-velocity (CV) target. Call a finite
control sequence $m$ a \emph{certified straight-tail escape} for state $s$ if its
tail segment is \emph{constant velocity and straight} ($\omega=0\wedge a=0$) and the
sound Lipschitz clearance certificate (the same interval lower bound used for
Prop.~\ref{prop:farfield}, Section~\ref{sec:provable-farfield}) shows hull
distance $d>0$ on $[0,H]$ with a past-closest-point-of-approach (CPA) margin. Define
the \emph{clearable set} $A:=\{s:\exists\,\text{certified straight-tail escape}\}$.

\begin{lemma}[Constant-velocity tail $\Rightarrow$ permanent clearance]
\label{lem:convex-tail}
If the tail is constant velocity and straight and the target is CV, the hull distance
$g(t)=\mathrm{dist}(\text{two fixed-orientation rectangles translating at constant
relative velocity})$ is \emph{convex} in $t$; hence once $g$ is increasing past the
CPA it increases for all later $t$, so $d>0$ on $[0,H]$ together with the past-CPA
margin implies $d>0$ for all $t\ge0$.
\end{lemma}
\emph{Remark (soundness fix).} Convexity requires the tail to be constant
\emph{velocity}: $\omega=0$ alone is insufficient, since an accelerating straight tail
traces a path that is non-affine in $t$ and $g$ need not be convex. The certificate
therefore requires $a=0$ as well. [HONESTY] The certificate is a \emph{sufficient}
condition; it is sound (never certifies a colliding escape---verified by a
long-horizon fuzz with zero false positives) but conservative (it rejects some safe
escapes).

\begin{proposition}[Control invariance of the clearable set]
\label{prop:forward-invariant}
For every $s\in A$ with certified escape $m=(u_0,u_1,\dots)$, applying the first
control $u_0$ over one decision step (official semantics, with the end-of-step speed
clamp aligned to $\Delta$) yields a successor $s'=F(s,u_0)$ whose tail
$m'=(u_1,u_2,\dots)$ is a certified straight-tail escape for $s'$ (it is the second
half of the \emph{same} physical trajectory, which passed one sound certificate).
Hence $s'\in A$: $A$ is control invariant. Consequently a terminal constraint
``successor $\in A$'', realized as an $O(1)$ re-certification of one concrete backup
maneuver (the backup-maneuver terminal set $U_{\mathrm{term}}$), renders the shielded
closed loop forward invariant on $A$, i.e.\ \emph{provably collision-free for all
future time on $A$ under a single CV target}.
\end{proposition}
\emph{Proof sketch.} The tail is certified because it is a sub-arc of an already
certified trajectory (verified numerically to reproduce the original trajectory's
second half exactly under the aligned clamp). Feasibility of $U_{\mathrm{term}}$ is
non-vacuous: $u_0$ itself keeps the backup certified, so the projection QP always
retains at least the backup as a fallback. $\square$

\paragraph{Compliant forward invariance ($A\cap U_{\mathrm{colregs}}$).}
The shield must pick a give-way-\emph{compliant} first step, so the relevant invariant
is $A\cap U_{\mathrm{colregs}}$. Let the compliant backup be ``turn to the mandated
side for $t_1$\,s, then go straight.'' \textbf{Case 1} ($t_1>\Delta$): the tail still
begins with a compliant turn and remains certified by
Prop.~\ref{prop:forward-invariant}, so $s'\in A\cap U_{\mathrm{colregs}}$---this case
is proved. \textbf{Case 2} ($t_1\le\Delta$, the turn completes within one step): the
tail's first step is straight and a fresh compliant escape from $s'$ is required; a
general proof of this case is open. On the shield's real give-way state distribution,
however, Case 2 \emph{never occurs} (the slow yaw rate, $\sim\!1.7^\circ$/s for a
$175$--$260$\,m hull, forces every compliant clearing turn to exceed one decision
step): across all sampled real give-way states with a compliant backup, $100\%$ fall
in Case 1, the tail remains certified and compliant in $100\%$, and the successor
re-admits a fresh compliant certified escape in $100\%$. We therefore state:
$A\cap U_{\mathrm{colregs}}$ is forward invariant via the proved Case 1 whenever the
compliant backup's turn extends beyond one decision step, a condition met on $100\%$
of the observed give-way distribution.

\begin{scope}
[HONESTY] (i) Single CV target (multi-target is weaker, $A=\bigcap_i A_i$, and per-%
target funnels may conflict). (ii) Membership in $A$ is decided by a finite
constant-control maneuver family plus the sound-but-conservative certificate, so $A$
is an \emph{inner} approximation of the true clearable set; the residual is exactly
the $A$-membership boundary---states not clearable in \emph{any} direction---which
falls to the emergency/relaxed fallback, and is \textbf{not} a
compliance-versus-safety conflict. On the real give-way distribution this residual is
$\approx\!0$--$3\%$ (head-on/overtake $0\%$, crossing $\sim\!2.9\%$), and every
in-$A$ give-way state admits a compliant certified escape. [HONESTY] These rates were
computed with a high-accuracy local integrator that matches the official dynamics to
sub-micron; official-pipeline reconfirmation is [PENDING server, evaluation-only].
(iii) The backup-maneuver terminal constraint $U_{\mathrm{term}}$ is implemented and
unit-tested as a sound module; wiring it into the deployed shield and validating the
closed loop is [PENDING server, evaluation-only]. (iv) The give-way states are drawn
from a well-behaved shielded rollout in which holding course never collides within the
horizon, so these rates characterize the clearable-set geometry, not performance under
adversarially imminent conflict, which needs a give-way-rich adversarial source
[PENDING].
\end{scope}

\subsection{Robust extension to multiple targets}
\label{sec:provable-multi}

For $N$ targets with occupancies $O_i(t)$, collision freedom is a conjunction over
targets, $\text{no collision}\iff\bigwedge_{i=1}^{N}(\text{ego\_body}\cap O_i=\varnothing)$.
This structure is the mathematical root of the robust extension: sufficient
conditions compose as $\forall i$, while the unavoidability certificate composes as
$\exists i$.

\begin{proposition}[Far-field, multi-target]
\label{prop:farfield-N}
If $d_i\ge d_{\mathrm{safe}}+v_{\mathrm{obs,max}}\Delta t+\rho_{\mathrm{ego}}$ for
\emph{all} $i$ (the same $\sim\!780$~m threshold), then after one step the ego body
intersects no target body.
\end{proposition}
The single-target proof holds independently for each $i$; the conjunction is exactly
the definition of collision freedom, so the extension is \emph{sound and tight}
(no extra conservatism beyond the per-pair circumscribed-disk bound). The new
obligation is (A2-N): $v_{\mathrm{obs}}^i\le9.5$ for \emph{each} target---the
single-target empirical bound must be re-verified on multi-target scenarios.

\begin{proposition}[Give-way, multi-target]
\label{prop:compliance-N}
Let $G=\{i:\rho_i \text{ is give-way}\}$ with compliant half-planes $H_i$. At every
step where $(i)$ $H_G:=\bigcap_{i\in G}H_i\ne\varnothing$ (the give-way directions
coexist) and $(ii)$ the QP over $U_{\mathrm{box}}\cap H_G\cap\bigcap_i
U_{\mathrm{collision\text{-}free}}^i$ is feasible, the executed $u_{\mathrm{safe}}$ is
direction-compliant for all $i\in G$.
\end{proposition}
The half-planes enter one QP and the projected solution satisfies their
intersection. The honest scope is sharper than in the single-target case: if two
targets demand opposite give-way directions then $H_G=\varnothing$ and the step falls
to fallback (a COLREGs \emph{violation}). This is not a defect but a boundary of
COLREGs itself---a pairwise rule set with no crisp multi-vessel norm (Rule~2 /
special-circumstances); our defensible position is ``compliant when the pairwise
norms coexist, collision-avoidance-first otherwise.'' As $N$ grows the feasible set
shrinks and the give-way compliance rate \emph{decreases}; this must be measured,
not assumed.

\begin{proposition}[Imminent unavoidability, multi-target]
\label{prop:imminent-N}
With $R_{\mathrm{box}}(t)$ as in Proposition~\ref{prop:imminent} (target-independent),
if $\exists\,i,\ \exists\,t^\ast$ ($\omega_{\max}t^\ast\le\pi/2$) with
$R_{\mathrm{box}}(t^\ast)\subseteq O_i(t^\ast)\oplus\mathrm{disk}(12.7)$, then the
collision is unavoidable.
\end{proposition}
The per-pair certificate for target $i$ implies every reachable trajectory collides
with $O_i$, hence with some target; the existential $\exists i$ is exactly what the
emergency layer needs, at cost $O(N\cdot|T|)$. The half that does \emph{not} compose
is avoidability: per-pair no-fire for all $i$ does \emph{not} imply global
avoidability (evading $A$ may force the ego into the future position of $B$), so
in the multi-target world no-fire is even weaker than for a single target.

\subsection{Honest positioning of the provable layer}
\label{sec:provable-summary}

Table~\ref{tab:provable} summarizes the guarantees, their status, and their scope.
The headline claim we can defend, with all assumptions attached, is: continuous
action $\cap$ provably \emph{direction}-compliant $\cap$ COLREGs (a previously empty
intersection), plus far-field single-step provable safety ($\forall i$, sound), a
sound imminent-unavoidability certificate for the emergency regime ($\exists i$), and
---the central upgrade---\emph{forward-invariant recursive feasibility} on a
certificate-defined clearable set (Prop.~\ref{prop:forward-invariant}): the shield is
provably collision-free for all future time on that set under a single CV target, a
guarantee reactive CBF-QP filters lack under a $10$\,s hold. We do \textbf{not} claim
bare provable collision freedom, complete provable coverage, maneuvering targets, or
multi-target direction conflicts (which fall to fallback by the pairwise nature of
COLREGs); the forward-invariance is scoped to the clearable set, whose residual
($\approx\!0$--$3\%$) is deferred to fallback. The provable tier is therefore
\emph{medium-to-strong}: the invariance upgrade and gap\#1 are sound, but coverage is
still an inner approximation, the closed-loop $U_{\mathrm{term}}$ wiring and official
reconfirmation are [PENDING server, eval-only], and the empirical zero-collision
backing remains single-target until multi-vessel benchmark data are available.

\begin{table}[h]
  \caption{Provable safety layer: guarantees, status, and honest scope}
  \label{tab:provable}
  \centering
  \small
  \begin{tabular}{@{}p{0.20\linewidth}p{0.24\linewidth}p{0.46\linewidth}@{}}
    \toprule
    \textbf{Guarantee} & \textbf{Status} & \textbf{Scope / limits} \\
    \midrule
    Far-field single-step no-collision (Prop.~\ref{prop:farfield},
    \ref{prop:farfield-N})
      & Proven (elementary; single-target bound conservative, multi-target conjunction adds no extra slack)
      & $d\ge\sim\!780$~m; CV target(s); $\forall i$ conjunction; each
        $v_{\mathrm{obs}}\le9.5$; single step, not invariance; early exit not yet
        implemented. \\
    \addlinespace
    Per-give-way-step direction compliance (Prop.~\ref{prop:compliance},
    \ref{prop:compliance-N})
      & By construction (scope-limited)
      & Give-way encounters ($\sim\!13$--$16\%$, encounter-level, not per step);
        excludes emergency $\rho_5$ and fallbacks; multi-target rate drops with $N$;
        direction conflict $\to$ fallback (COLREGs pairwise boundary). Never ``every
        step.'' \\
    \addlinespace
    Imminent-unavoidability certificate (gap\#1) (Prop.~\ref{prop:imminent},
    \ref{prop:imminent-N})
      & Sound sufficient condition ($>\!20{,}000$ adversarial scenarios, 0 false positives)
      & Imminent only; CV target(s); $\exists i$ composes, avoidability does not;
        not complete / not early-warning; a classifier, not a system guarantee. \\
    \addlinespace
    Forward-invariant recursive feasibility on the clearable set $A$
    (Lem.~\ref{lem:convex-tail}, Prop.~\ref{prop:forward-invariant})
      & Proven (control invariance of $A$; certificate sound, 0 false positives on
        long-horizon fuzz). $U_{\mathrm{term}}$ module unit-tested; closed-loop wiring
        [PENDING server, eval-only]
      & Single CV target; $A$ = inner approximation (finite maneuver family +
        conservative certificate); residual $\approx\!0$--$3\%$ = states clearable in
        no direction $\to$ fallback (not a compliance conflict); $A\cap U_{\mathrm{colregs}}$
        invariance proved for Case 1 and holds on $100\%$ of the observed give-way
        distribution; numbers local-integrator, official reconfirmation [PENDING]. \\
    \bottomrule
  \end{tabular}
\end{table}

\paragraph{Emergency fallback and the failure-mechanism boundary.}
When the projection is infeasible ($P=\varnothing$) in the emergency regime
$\rho_5$, the \texttt{is\_emergency} predicate structurally bypasses the direction
constraint and Algorithm~1 minimizes collision risk---an \emph{empirical} fallback,
not a guarantee. On an adversarial purely-geometric baseline that forces the shield
into the emergency branch on $44$--$69\%$ of steps, we observe $45\%$ arrival and
two collisions; auditing confirms both occur with
\texttt{source=emergency},~$\rho=5$, i.e.\ the emergency last resort trying and
failing to avoid, and both lie outside the domain of Proposition~\ref{prop:farfield}
(far-field, non-emergency, single-step). They are not counterexamples to the headline
but they do substantiate that the emergency regime is not a guarantee. On the learned
policy (non-adversarial), we observe zero hull collisions over $10$ seeds
$\times$ $40$ episodes ($400$ total), so the ``empirically zero-collision'' statement holds; we
nonetheless \textbf{never} write provably collision-free and label the emergency
regime as an empirical fallback throughout.
 进 Methods 之后（或作为独立 section）。
%  依赖宏包：\usepackage{amsmath,amsthm,amssymb,bm}。TRB 无附录/7500 词 →
%  正文留命题陈述 + 证明梗概 + 诚实作用域；完整证明移交外部技术报告（见 \cite{techreport}）。
% =====================================================================

% ---- 若主文档尚未定义，取消注释以下 theorem 环境 ----
% \newtheorem{proposition}{Proposition}
% \newtheorem{assumption}{Assumption}
% \theoremstyle{remark}\newtheorem*{scope}{Scope and honest limits}

\section{The Provable Safety Layer}
\label{sec:provable}

Our contribution is \emph{not} a claim of provable collision freedom. Prior
provably COLREGs-\emph{direction}-compliant collision avoidance either sacrificed
the continuous action space (discrete action masking, Krasowski and Althoff 2024)
or sacrificed the guarantee (continuous-action work reverts to soft reward
shaping). We place a \emph{projection-based safety shield} on maritime COLREGs in
the continuous action space and characterize \emph{exactly what can and cannot be
proven}. The safety layer is organized by encounter regime and yields three
guarantees of different strengths: a far-field single-step no-collision guarantee
(Proposition~\ref{prop:farfield}), a per-give-way-step direction-compliance
guarantee (Proposition~\ref{prop:compliance}), and a \emph{sound} imminent-%
unavoidability certificate for the emergency regime (Proposition~\ref{prop:imminent}).
Each is stated with its full assumption set, and each carries an explicit scope
boundary. Full proofs are deferred to the external technical report; here we give
statements and proof sketches. Throughout, ``proven'' means a strict lemma;
existence-level and empirical results are labeled as such.

\subsection{System setting and notation}
\label{sec:provable-setup}

We adopt the yaw-constrained point-mass vessel dynamics of Krasowski and Althoff
(2024) (vessel class SR108 container ship, $175 \times 25.4$~m):
\begin{equation}
  \dot{x}=v\cos\theta,\quad \dot{y}=v\sin\theta,\quad \dot{\theta}=\omega,\quad \dot{v}=a,
  \label{eq:dyn}
\end{equation}
with control bounds $|a|\le a_{\max}=0.24~\mathrm{m/s^2}$,
$|\omega|\le\omega_{\max}=0.03~\mathrm{rad/s}$, speed clipped by the environment to
$0\le v\le v_{\max}=9.5~\mathrm{m/s}$ each step (within a step $v$ may briefly reach
$v_{\max}+a_{\max}\Delta t=11.9~\mathrm{m/s}$), and decision step
$\Delta t=10~\mathrm{s}$. The ego body is the $175\times25.4$~m rectangle with
circumscribed radius $R_{\mathrm{circ}}=\tfrac12\sqrt{175^2+25.4^2}=88.4$~m and
inscribed radius $r_{\mathrm{insc}}=\tfrac12\cdot25.4=12.7$~m. Target vessels occupy
their body rectangle (or, conservatively, their circumscribed disk); unless stated
otherwise, in stand-on regimes they are extrapolated at \emph{constant velocity}
(CV). We write $p_e,p_m$ for ego/target positions,
$\bm{v}_e=v_e(\cos\theta_e,\sin\theta_e)$ for the ego velocity vector, and
$\hat{h},\hat{h}^{\perp}$ for the heading unit vector and its normal.

\subsection{Far-field regime: single-step no-collision}
\label{sec:provable-farfield}

\begin{proposition}[Far-field single-step no-collision]
\label{prop:farfield}
Let $d=\lVert p_e-p_m\rVert$ be the ego--target center distance for a single target.
If
\begin{equation}
  d \;\ge\; d_{\mathrm{safe}} + v_{\mathrm{obs,max}}\,\Delta t + \rho_{\mathrm{ego}},
  \label{eq:farfield}
\end{equation}
then after one decision step the two (circumscribed-disk) bodies cannot intersect,
i.e.\ a collision is impossible in that step. Here
$d_{\mathrm{safe}}=R_{\mathrm{circ,ego}}+R_{\mathrm{circ,obs}}$ and
$\rho_{\mathrm{ego}}=v_{\max}\Delta t+\tfrac12 a_{\max}\Delta t^2$ is the maximum
single-step ego displacement.
\end{proposition}

Substituting constants gives $R_{\mathrm{circ,ego}}=88.4$~m and, because the target
width is not in the state, a conservative square upper bound with the Table~II
safety margin $d_{\mathrm{obs,safety}}=2\ell$:
$R_{\mathrm{circ,obs}}=\tfrac12\sqrt2\,\ell+d_{\mathrm{obs,safety}}
=123.74+350=473.74$~m, so $d_{\mathrm{safe}}=562.16$~m,
$\rho_{\mathrm{ego}}=9.5\cdot10+\tfrac12\cdot0.24\cdot100=107.0$~m, and
$v_{\mathrm{obs,max}}\Delta t=95$~m, yielding a threshold of $764.16$~m
(engineering value $\sim\!780$~m with margin).

\emph{Assumptions.} (A1) single target (for $N$ targets, each $i$ must satisfy the
threshold, Proposition~\ref{prop:farfield-N}); (A2) target speed
$v_{\mathrm{obs}}\le9.5~\mathrm{m/s}$---Krasowski sets $v_{\max}$ only for the ego, so
this bound must be assumed explicitly. It is \emph{empirically verified} on the full
benchmark: across all $2000$ HandcraftedTwoVessel scenarios ($342{,}000$ obstacle
states) the maximum is $7.10~\mathrm{m/s}$, with zero states exceeding $9.5$
($\sim\!2.4~\mathrm{m/s}$ margin); higher-speed real AIS data would require raising
the threshold. (A3) $v\le v_{\max}$ at step end (environment clip); (A4) the target
is approximately CV within a step (stand-on); (A5) circumscribed-disk occupancy
(disks disjoint $\Rightarrow$ bodies disjoint).

\emph{Proof sketch.} By the triangle inequality the next-step center distance
satisfies $d_{\mathrm{next}}\ge d-\lVert\Delta p_e\rVert-\lVert\Delta p_m\rVert
\ge d-\rho_{\mathrm{ego}}-v_{\mathrm{obs,max}}\Delta t$. With \eqref{eq:farfield},
$d_{\mathrm{next}}\ge d_{\mathrm{safe}}=R_{\mathrm{circ,ego}}+R_{\mathrm{circ,obs}}$,
so the circumscribed disks---hence the bodies---are disjoint. $\square$

\begin{scope}
The triangle inequality is direction-independent, so head-on is already the true
worst case and there is no geometric gap; this is the \emph{soundest} block of the
provable layer (an elementary lemma). It is a conservative single-step
\emph{sufficient} condition---not invariance, not a near-field guarantee---and is
currently not implemented as an early exit (an $O(1)$ far-field fast path could be
added).
\end{scope}

\subsection{Give-way regime: per-step direction compliance}
\label{sec:provable-compliance}

\begin{proposition}[Per-give-way-step direction compliance]
\label{prop:compliance}
At every decision step that the state machine classifies as a give-way regime
($\rho\in\{\text{head-on},\text{crossing},\text{overtake}\}$) and for which the
projection quadratic program (QP) is feasible (\texttt{needs\_fallback}=False), the
executed safe action $u_{\mathrm{safe}}$ lies in the COLREGs-mandated compliant
half-plane (e.g.\ give-way starboard turn $\omega\le-\omega_{\mathrm{turn}}$).
\end{proposition}

\emph{Assumptions.} (B1) single target; (B2) correct state-machine classification of
$\rho$ (inheriting the $\sim\!12.6\%$ classification gap and the Lemma-1 tie-break
issue); (B3) QP feasibility (not fallback); (B4) scope $=$ give-way encounters. The
empirical give-way trigger rate is $\sim\!13$--$16\%$ \emph{at the
encounter/episode level} (measured twice over $933$ scenarios; overtake $=0$)---not
per step; stand-on violations are counted per step.

\emph{Proof.} The compliant half-plane enters the QP as a hard interval constraint
($w_{\mathrm{hi}}=-\omega_{\mathrm{turn}}$, etc.), so any projected solution
satisfies it. $\square$ (by construction; no separate theorem required).

\begin{scope}
This does \emph{not} cover the emergency regime $\rho_5$ (which occupies
$\sim\!55$--$60\%$ of conflict steps and whose \texttt{is\_emergency} predicate
structurally bypasses the compliance constraint), nor the relaxed/%
collision-minimizing fallbacks (which drop the half-plane). We therefore write only
``provably give-way-direction-compliant, per give-way step'' and \textbf{never}
``provably compliant / compliant at every step'' (the measured violation rate on the
golden 10-seed run is a nonzero $1.13$--$1.8$ per episode).
\end{scope}

\subsection{Emergency regime: a sound imminent-unavoidability certificate}
\label{sec:provable-imminent}

The emergency controller is an empirical last resort. To characterize it honestly we
do not claim it avoids collisions; instead we provide a \emph{sound sufficient
condition} that certifies when a collision is physically unavoidable, so that the
``black box'' becomes a decidable boundary. This replaces an earlier
turn-only condition that was \emph{unsound} because it ignored acceleration/%
deceleration escapes (gap\#1).

\begin{proposition}[Imminent-unavoidability certificate, single CV target]
\label{prop:imminent}
Consider a single constant-velocity (stand-on) target. Define the ego
reachable-center over-approximation $R_{\mathrm{box}}(t)$ (a heading-frame rectangle
centered at the CV extrapolation $c_e(t)=p_e+\bm{v}_e t$) with longitudinal
half-axis $B_{\mathrm{lon}}(t)=\tfrac12 a_{\mathrm{eff}}t^2$ and lateral half-axis
$L_{\mathrm{lat}}(t)$, where
\begin{align}
  a_{\mathrm{eff}} &= \sqrt{a_{\max}^2+(v_{\mathrm{bnd}}\,\omega_{\max})^2},
    \quad v_{\mathrm{bnd}}=v_{\max}+a_{\max}\Delta t=11.9~\mathrm{m/s},\\
  L_{\mathrm{lat}}(t) &=
    \begin{cases}
      \dfrac{v_{\mathrm{bnd}}}{\omega_{\max}}\bigl(1-\cos\omega_{\max}t\bigr),
        & \omega_{\max}t\le\tfrac{\pi}{2},\\[1.2ex]
      \dfrac{v_{\mathrm{bnd}}}{\omega_{\max}}
        + v_{\mathrm{bnd}}\Bigl(t-\tfrac{\pi}{2\omega_{\max}}\Bigr),
        & \text{otherwise.}
    \end{cases}
\end{align}
Let $O(t)$ be the target body rectangle (center $p_m+\bm{v}_m t$, heading
$\theta_m$). If
\begin{equation}
  \exists\, t^\ast\in[0,T]:\quad
  R_{\mathrm{box}}(t^\ast)\;\subseteq\;\bigl(O(t^\ast)\oplus \mathrm{disk}(r_{\mathrm{insc,ego}}=12.7~\mathrm{m})\bigr),
  \label{eq:imminent}
\end{equation}
then the collision is \emph{unavoidable}: every feasible control sequence collides
with the target at $t^\ast$.
\end{proposition}

\emph{Assumptions.} (A1) single target, \emph{per-pair} (per-pair unavoidable
$\Rightarrow$ global unavoidable; per-pair \emph{avoidable} does \emph{not}
compose); (A2) target CV stand-on (a maneuvering target changes $O(t)$ and voids the
guarantee---we state ``unavoidable within the CV world''); (A3) $v\le
v_{\mathrm{bnd}}=11.9$ within a step (includes the clip overshoot); (A4) the ego
inscribed disk $r_{\mathrm{insc}}=12.7~\mathrm{m}\subseteq$ body at any heading
(from the short edge; the long-edge zero-margin case is pending an independent re-check); (A5)
the certificate fires only in the region $\omega_{\max}t^\ast\le\pi/2$
($\le52.4$~s)---we verified that firing is limited by the longitudinal quadratic term
to $t_{\mathrm{crit,lon}}=\sqrt{2(\tfrac12\cdot175+12.7)/a_{\mathrm{eff}}}=21.6~\mathrm{s}<52.4~\mathrm{s}$;
(A6) single-instant containment (sufficient, not necessary); (A7) occupancy
under-approximation: $\text{obs\_width}\le$ true target width and $L_{\mathrm{obs}}\le$
true target length (a soundness precondition the caller must guarantee---passing
sizes larger than the truth over-approximates $O$ and can yield false positives).

\emph{Proof sketch (three steps).} (i) $R_{\mathrm{box}}$ over-approximates the true
reachable center: $\lVert d\bm{v}/dt\rVert=\sqrt{a^2+(v\omega)^2}\le a_{\mathrm{eff}}$,
and double integration bounds the longitudinal offset by $\tfrac12 a_{\mathrm{eff}}t^2$
while $|y'(t)|\le\int_0^t v_{\mathrm{bnd}}\sin(\omega_{\max}s)\,ds=L_{\mathrm{lat}}(t)$.
(ii) The inscribed disk under-approximates the body. (iii) If
$R_{\mathrm{box}}(t^\ast)\subseteq O(t^\ast)\oplus\mathrm{disk}(12.7)$, then every
reachable ego center lies within $12.7$~m of $O(t^\ast)$, so the ego inscribed disk
($\subseteq$ body) intersects $O(t^\ast)$: a collision. Since this holds for
\emph{all} reachable positions, every trajectory collides at $t^\ast$. $\square$

\emph{Verification.} A six-agent adversarial audit ($20{,}000$ random $+$ $735$
curated cases, with a triple-switch bang-bang escape search under exact hull-collision
geometry) found \emph{zero} true false positives; after hardening the
$v_{\mathrm{bnd}}/L_{\mathrm{lat}}$ envelope, a battery plus $1500$ fuzz cases
remained sound (0 false positives) and non-vacuous (firing at $t^\ast\le7$~s).

\begin{scope}
The certificate is a conservative \emph{one-sided sufficient} condition for
imminent unavoidability. It is \textbf{not} an early-warning cost (it is a terminal
classifier, not a planner), \textbf{not} complete or necessary (no-fire $\ne$
avoidable), and by itself certifies nothing about collision freedom. We
\textbf{never} write ``provably zero-collision'': this is an emergency-regime
classifier, not a system-wide guarantee.
\end{scope}

\subsection{Forward-invariant recursive feasibility: the provable upgrade}
\label{sec:provable-forward-invariant}

Propositions~\ref{prop:farfield}--\ref{prop:compliance} are \emph{one-step}
guarantees; a single-step no-collision constraint is not by itself recursively
feasible (a state may be safe now yet admit no safe continuation). This subsection
upgrades the shield from one-step lookahead to a \emph{forward-invariant} guarantee
on a certified clearable set, which is the paper's central provable contribution and
the differentiator against reactive CBF-QP filters (which, under a $\Delta=10$\,s
zero-order hold, carry no recursive-feasibility proof).

\paragraph{Clearable set.} Fix a single constant-velocity (CV) target. Call a finite
control sequence $m$ a \emph{certified straight-tail escape} for state $s$ if its
tail segment is \emph{constant velocity and straight} ($\omega=0\wedge a=0$) and the
sound Lipschitz clearance certificate (the same interval lower bound used for
Prop.~\ref{prop:farfield}, Section~\ref{sec:provable-farfield}) shows hull
distance $d>0$ on $[0,H]$ with a past-closest-point-of-approach (CPA) margin. Define
the \emph{clearable set} $A:=\{s:\exists\,\text{certified straight-tail escape}\}$.

\begin{lemma}[Constant-velocity tail $\Rightarrow$ permanent clearance]
\label{lem:convex-tail}
If the tail is constant velocity and straight and the target is CV, the hull distance
$g(t)=\mathrm{dist}(\text{two fixed-orientation rectangles translating at constant
relative velocity})$ is \emph{convex} in $t$; hence once $g$ is increasing past the
CPA it increases for all later $t$, so $d>0$ on $[0,H]$ together with the past-CPA
margin implies $d>0$ for all $t\ge0$.
\end{lemma}
\emph{Remark (soundness fix).} Convexity requires the tail to be constant
\emph{velocity}: $\omega=0$ alone is insufficient, since an accelerating straight tail
traces a path that is non-affine in $t$ and $g$ need not be convex. The certificate
therefore requires $a=0$ as well. [HONESTY] The certificate is a \emph{sufficient}
condition; it is sound (never certifies a colliding escape---verified by a
long-horizon fuzz with zero false positives) but conservative (it rejects some safe
escapes).

\begin{proposition}[Control invariance of the clearable set]
\label{prop:forward-invariant}
For every $s\in A$ with certified escape $m=(u_0,u_1,\dots)$, applying the first
control $u_0$ over one decision step (official semantics, with the end-of-step speed
clamp aligned to $\Delta$) yields a successor $s'=F(s,u_0)$ whose tail
$m'=(u_1,u_2,\dots)$ is a certified straight-tail escape for $s'$ (it is the second
half of the \emph{same} physical trajectory, which passed one sound certificate).
Hence $s'\in A$: $A$ is control invariant. Consequently a terminal constraint
``successor $\in A$'', realized as an $O(1)$ re-certification of one concrete backup
maneuver (the backup-maneuver terminal set $U_{\mathrm{term}}$), renders the shielded
closed loop forward invariant on $A$, i.e.\ \emph{provably collision-free for all
future time on $A$ under a single CV target}.
\end{proposition}
\emph{Proof sketch.} The tail is certified because it is a sub-arc of an already
certified trajectory (verified numerically to reproduce the original trajectory's
second half exactly under the aligned clamp). Feasibility of $U_{\mathrm{term}}$ is
non-vacuous: $u_0$ itself keeps the backup certified, so the projection QP always
retains at least the backup as a fallback. $\square$

\paragraph{Compliant forward invariance ($A\cap U_{\mathrm{colregs}}$).}
The shield must pick a give-way-\emph{compliant} first step, so the relevant invariant
is $A\cap U_{\mathrm{colregs}}$. Let the compliant backup be ``turn to the mandated
side for $t_1$\,s, then go straight.'' \textbf{Case 1} ($t_1>\Delta$): the tail still
begins with a compliant turn and remains certified by
Prop.~\ref{prop:forward-invariant}, so $s'\in A\cap U_{\mathrm{colregs}}$---this case
is proved. \textbf{Case 2} ($t_1\le\Delta$, the turn completes within one step): the
tail's first step is straight and a fresh compliant escape from $s'$ is required; a
general proof of this case is open. On the shield's real give-way state distribution,
however, Case 2 \emph{never occurs} (the slow yaw rate, $\sim\!1.7^\circ$/s for a
$175$--$260$\,m hull, forces every compliant clearing turn to exceed one decision
step): across all sampled real give-way states with a compliant backup, $100\%$ fall
in Case 1, the tail remains certified and compliant in $100\%$, and the successor
re-admits a fresh compliant certified escape in $100\%$. We therefore state:
$A\cap U_{\mathrm{colregs}}$ is forward invariant via the proved Case 1 whenever the
compliant backup's turn extends beyond one decision step, a condition met on $100\%$
of the observed give-way distribution.

\begin{scope}
[HONESTY] (i) Single CV target (multi-target is weaker, $A=\bigcap_i A_i$, and per-%
target funnels may conflict). (ii) Membership in $A$ is decided by a finite
constant-control maneuver family plus the sound-but-conservative certificate, so $A$
is an \emph{inner} approximation of the true clearable set; the residual is exactly
the $A$-membership boundary---states not clearable in \emph{any} direction---which
falls to the emergency/relaxed fallback, and is \textbf{not} a
compliance-versus-safety conflict. On the real give-way distribution this residual is
$\approx\!0$--$3\%$ (head-on/overtake $0\%$, crossing $\sim\!2.9\%$), and every
in-$A$ give-way state admits a compliant certified escape. [HONESTY] These rates were
computed with a high-accuracy local integrator that matches the official dynamics to
sub-micron; official-pipeline reconfirmation is [PENDING server, evaluation-only].
(iii) The backup-maneuver terminal constraint $U_{\mathrm{term}}$ is implemented and
unit-tested as a sound module; wiring it into the deployed shield and validating the
closed loop is [PENDING server, evaluation-only]. (iv) The give-way states are drawn
from a well-behaved shielded rollout in which holding course never collides within the
horizon, so these rates characterize the clearable-set geometry, not performance under
adversarially imminent conflict, which needs a give-way-rich adversarial source
[PENDING].
\end{scope}

\subsection{Robust extension to multiple targets}
\label{sec:provable-multi}

For $N$ targets with occupancies $O_i(t)$, collision freedom is a conjunction over
targets, $\text{no collision}\iff\bigwedge_{i=1}^{N}(\text{ego\_body}\cap O_i=\varnothing)$.
This structure is the mathematical root of the robust extension: sufficient
conditions compose as $\forall i$, while the unavoidability certificate composes as
$\exists i$.

\begin{proposition}[Far-field, multi-target]
\label{prop:farfield-N}
If $d_i\ge d_{\mathrm{safe}}+v_{\mathrm{obs,max}}\Delta t+\rho_{\mathrm{ego}}$ for
\emph{all} $i$ (the same $\sim\!780$~m threshold), then after one step the ego body
intersects no target body.
\end{proposition}
The single-target proof holds independently for each $i$; the conjunction is exactly
the definition of collision freedom, so the extension is \emph{sound and tight}
(no extra conservatism beyond the per-pair circumscribed-disk bound). The new
obligation is (A2-N): $v_{\mathrm{obs}}^i\le9.5$ for \emph{each} target---the
single-target empirical bound must be re-verified on multi-target scenarios.

\begin{proposition}[Give-way, multi-target]
\label{prop:compliance-N}
Let $G=\{i:\rho_i \text{ is give-way}\}$ with compliant half-planes $H_i$. At every
step where $(i)$ $H_G:=\bigcap_{i\in G}H_i\ne\varnothing$ (the give-way directions
coexist) and $(ii)$ the QP over $U_{\mathrm{box}}\cap H_G\cap\bigcap_i
U_{\mathrm{collision\text{-}free}}^i$ is feasible, the executed $u_{\mathrm{safe}}$ is
direction-compliant for all $i\in G$.
\end{proposition}
The half-planes enter one QP and the projected solution satisfies their
intersection. The honest scope is sharper than in the single-target case: if two
targets demand opposite give-way directions then $H_G=\varnothing$ and the step falls
to fallback (a COLREGs \emph{violation}). This is not a defect but a boundary of
COLREGs itself---a pairwise rule set with no crisp multi-vessel norm (Rule~2 /
special-circumstances); our defensible position is ``compliant when the pairwise
norms coexist, collision-avoidance-first otherwise.'' As $N$ grows the feasible set
shrinks and the give-way compliance rate \emph{decreases}; this must be measured,
not assumed.

\begin{proposition}[Imminent unavoidability, multi-target]
\label{prop:imminent-N}
With $R_{\mathrm{box}}(t)$ as in Proposition~\ref{prop:imminent} (target-independent),
if $\exists\,i,\ \exists\,t^\ast$ ($\omega_{\max}t^\ast\le\pi/2$) with
$R_{\mathrm{box}}(t^\ast)\subseteq O_i(t^\ast)\oplus\mathrm{disk}(12.7)$, then the
collision is unavoidable.
\end{proposition}
The per-pair certificate for target $i$ implies every reachable trajectory collides
with $O_i$, hence with some target; the existential $\exists i$ is exactly what the
emergency layer needs, at cost $O(N\cdot|T|)$. The half that does \emph{not} compose
is avoidability: per-pair no-fire for all $i$ does \emph{not} imply global
avoidability (evading $A$ may force the ego into the future position of $B$), so
in the multi-target world no-fire is even weaker than for a single target.

\subsection{Honest positioning of the provable layer}
\label{sec:provable-summary}

Table~\ref{tab:provable} summarizes the guarantees, their status, and their scope.
The headline claim we can defend, with all assumptions attached, is: continuous
action $\cap$ provably \emph{direction}-compliant $\cap$ COLREGs (a previously empty
intersection), plus far-field single-step provable safety ($\forall i$, sound), a
sound imminent-unavoidability certificate for the emergency regime ($\exists i$), and
---the central upgrade---\emph{forward-invariant recursive feasibility} on a
certificate-defined clearable set (Prop.~\ref{prop:forward-invariant}): the shield is
provably collision-free for all future time on that set under a single CV target, a
guarantee reactive CBF-QP filters lack under a $10$\,s hold. We do \textbf{not} claim
bare provable collision freedom, complete provable coverage, maneuvering targets, or
multi-target direction conflicts (which fall to fallback by the pairwise nature of
COLREGs); the forward-invariance is scoped to the clearable set, whose residual
($\approx\!0$--$3\%$) is deferred to fallback. The provable tier is therefore
\emph{medium-to-strong}: the invariance upgrade and gap\#1 are sound, but coverage is
still an inner approximation, the closed-loop $U_{\mathrm{term}}$ wiring and official
reconfirmation are [PENDING server, eval-only], and the empirical zero-collision
backing remains single-target until multi-vessel benchmark data are available.

\begin{table}[h]
  \caption{Provable safety layer: guarantees, status, and honest scope}
  \label{tab:provable}
  \centering
  \small
  \begin{tabular}{@{}p{0.20\linewidth}p{0.24\linewidth}p{0.46\linewidth}@{}}
    \toprule
    \textbf{Guarantee} & \textbf{Status} & \textbf{Scope / limits} \\
    \midrule
    Far-field single-step no-collision (Prop.~\ref{prop:farfield},
    \ref{prop:farfield-N})
      & Proven (elementary; single-target bound conservative, multi-target conjunction adds no extra slack)
      & $d\ge\sim\!780$~m; CV target(s); $\forall i$ conjunction; each
        $v_{\mathrm{obs}}\le9.5$; single step, not invariance; early exit not yet
        implemented. \\
    \addlinespace
    Per-give-way-step direction compliance (Prop.~\ref{prop:compliance},
    \ref{prop:compliance-N})
      & By construction (scope-limited)
      & Give-way encounters ($\sim\!13$--$16\%$, encounter-level, not per step);
        excludes emergency $\rho_5$ and fallbacks; multi-target rate drops with $N$;
        direction conflict $\to$ fallback (COLREGs pairwise boundary). Never ``every
        step.'' \\
    \addlinespace
    Imminent-unavoidability certificate (gap\#1) (Prop.~\ref{prop:imminent},
    \ref{prop:imminent-N})
      & Sound sufficient condition ($>\!20{,}000$ adversarial scenarios, 0 false positives)
      & Imminent only; CV target(s); $\exists i$ composes, avoidability does not;
        not complete / not early-warning; a classifier, not a system guarantee. \\
    \addlinespace
    Forward-invariant recursive feasibility on the clearable set $A$
    (Lem.~\ref{lem:convex-tail}, Prop.~\ref{prop:forward-invariant})
      & Proven (control invariance of $A$; certificate sound, 0 false positives on
        long-horizon fuzz). $U_{\mathrm{term}}$ module unit-tested; closed-loop wiring
        [PENDING server, eval-only]
      & Single CV target; $A$ = inner approximation (finite maneuver family +
        conservative certificate); residual $\approx\!0$--$3\%$ = states clearable in
        no direction $\to$ fallback (not a compliance conflict); $A\cap U_{\mathrm{colregs}}$
        invariance proved for Case 1 and holds on $100\%$ of the observed give-way
        distribution; numbers local-integrator, official reconfirmation [PENDING]. \\
    \bottomrule
  \end{tabular}
\end{table}

\paragraph{Emergency fallback and the failure-mechanism boundary.}
When the projection is infeasible ($P=\varnothing$) in the emergency regime
$\rho_5$, the \texttt{is\_emergency} predicate structurally bypasses the direction
constraint and Algorithm~1 minimizes collision risk---an \emph{empirical} fallback,
not a guarantee. On an adversarial purely-geometric baseline that forces the shield
into the emergency branch on $44$--$69\%$ of steps, we observe $45\%$ arrival and
two collisions; auditing confirms both occur with
\texttt{source=emergency},~$\rho=5$, i.e.\ the emergency last resort trying and
failing to avoid, and both lie outside the domain of Proposition~\ref{prop:farfield}
(far-field, non-emergency, single-step). They are not counterexamples to the headline
but they do substantiate that the emergency regime is not a guarantee. On the learned
policy (non-adversarial), we observe zero hull collisions over $10$ seeds
$\times$ $40$ episodes ($400$ total), so the ``empirically zero-collision'' statement holds; we
nonetheless \textbf{never} write provably collision-free and label the emergency
regime as an empirical fallback throughout.


% =====================================================================
\section{Experiments and Results}
\label{sec:results}
% [DRAFT skeleton · 大量 PENDING = 待烧卡]
\subsection{Setup}
Benchmark: Krasowski--Althoff CommonOcean handcrafted two-vessel scenarios
(head-on / crossing / overtake). Metrics (reported as IQM with bootstrap CI over
seeds, never cherry-picked): arrival rate, hull-collision rate, COLREGs violation
rate, and control smoothness (throttle-increment, yaw-increment, jerk). We report
two arrival criteria: strict (position + heading $\pm9.74^\circ$ + time, faithful
to the CommonOcean \texttt{is\_reached}) as primary, and position-only as a clearly
labeled looser complement. [HONESTY: dual criteria because Krasowski's published
maritime eval code is not in hand.]

\subsection{Main comparison (four arms)}
% [PENDING: headline four-way money chart on the balanced/overtaking dataset]
Baseline / soft-reward / discrete-safe / continuous-safe on the balanced dataset
(head-on + crossing + overtake). Expected honest positioning:
\begin{itemize}
  \item \textbf{Collision}: zero hull collisions for the shielded continuous arm
  (empirical; [PENDING seeds$\times$episodes]); safety methods tie near zero.
  \item \textbf{Throttle smoothness}: structural $\sim\!5\times$ advantage of the
  continuous arm over discrete (discrete has 7 quantized throttle levels).
  \textbf{[HONESTY] yaw-increment and jerk currently \emph{lose} to discrete
  ($\sim\!1.2$--$1.3\times$); de-jittering the yaw axis is open work.}
  \item \textbf{Arrival}: comparable to discrete, \emph{not} superior.
  [HONESTY: after balancing hidden collapsed seeds across arms, arrival shows no
  provable advantage over discrete (L192); we therefore do not list arrival as a
  selling point.]
  \item \textbf{Violation}: [PENDING same-budget same-dataset paired test.]
\end{itemize}

\subsection{Training-stability collapse and the warm-start remedy}
% [DRAFT · from 失败机理与紧急兜底刻画_0714 + L190/L191/L191-B/L192]
A minority of seeds (2/10 in the from-scratch golden run) collapse: the policy
never escapes a ``full-speed circling / time-out'' attractor and arrival flat-lines
near zero. We show (position-vs-heading decomposition; limit-cycle trajectories)
that this is a \emph{training-instability} limitation, not a criterion artifact,
and that reward shaping does not fix it (14 dosage sweeps). Warm-starting all seeds
from a single healthy policy escapes the bad basin (collapsed seeds' failure mode
changes entirely from time-out to near-goal stopping; zero hull collisions
maintained). [HONESTY: we describe this as ``warm-started training,'' never
``stable from scratch''; the ten warm-started seeds share one parent, so their
low spread is a fine-tuning band, \emph{not} a seed-variance claim---establishing
seed-level stability requires a multi-parent design, PENDING.]

\subsection{Emergency-regime characterization}
On an adversarial purely-geometric baseline that forces the emergency branch, we
observe the two collisions occurring with \texttt{source=emergency}, substantiating
that the emergency regime is an empirical fallback outside the far-field domain of
Proposition~\ref{prop:farfield}. [HONESTY: not a counterexample to the headline;
the headline is the provable direction-compliant shield with empirical zero
collision on the learned policy.]

% =====================================================================
\section{Discussion and Limitations}
% [DRAFT · 13 条 over-claim 收紧措辞已落此处 · 依据 03 L189 E + L191-B + L192]
% 数字来源：violation 1.13-1.8 / 崩率~20% CI[2.5,55.6] / 0-38 shared + 6.3-of-8 /
%   接管步 27-43% 路径~20% / 油门~5x 住 RL 段 / yaw 输 1.19-1.33x / 覆盖 96.32 vs 97.90.
\begin{itemize}
  \item \textbf{Provable tier is \emph{medium}.} gap\#1 (velocity-dimension
  unavoidability) is now sound, but coverage is not complete, and the empirical
  zero-collision backing is single-target until multi-vessel benchmark data exist.
  We keep every claim hedged and \emph{never} abbreviate the shield to ``provably
  collision-free.''

  \item \textbf{Compliance is per give-way step, not every step.} We claim provable
  \emph{direction} compliance only on give-way steps ($\sim\!13$--$16\%$ of
  encounters), never at every step (measured violation $1.13$--$1.8$ per episode on
  the golden run; the emergency regime structurally bypasses the constraint).

  \item \textbf{Arrival is comparable, not superior.} The apparent 24-point
  arrival win over discrete was a collapse-balancing artifact: after balancing the
  hidden collapsed seeds across arms, arrival shows \emph{no} statistically
  demonstrable advantage over discrete ($n=3$ per arm, sign-test $p=0.25$--$0.63$).
  We therefore do \emph{not} list arrival as a selling point.

  \item \textbf{The smoothness advantage is throttle-only and dilutes with
  takeover.} The $\sim\!5$--$6\times$ throttle-increment advantage is structural
  (discrete has 7 quantized throttle levels) but lives in the \emph{RL} segment; on
  the geometric-takeover segment the two arms are a wash, so the advantage shrinks
  as the takeover fraction grows. Yaw-increment and jerk currently \emph{lose} to
  discrete ($\sim\!1.19$--$1.33\times$); de-jittering the yaw axis is open work.

  \item \textbf{A minority-seed training-stability collapse.} A minority of seeds
  ($\sim\!20\%$, wide CI $[2.5,55.6]$) collapse. This is training instability, not
  scenario hardness: the shared-failure scenarios are $0/38$ unsolvable by all
  healthy seeds, and each is solved by $6.3/8$ healthy seeds on average (a clean
  seed-only control). Reward shaping is a \emph{high-confidence empirical} dead end
  for the tried shaping/coefficient families, \emph{not} a proof of impossibility
  for CMDP or non-PBRS terminal objectives.

  \item \textbf{Warm-start escapes the bad basin; it is not from-scratch
  stability.} All ten warm-started seeds derive from a single, post-hoc-selected
  parent, so their low spread is a fine-tuning band with effective sample size
  $\approx 1$, \emph{not} a seed-variance claim. Establishing seed-level stability
  requires the multi-parent design [PENDING]. Warm-start also slightly lowers
  provable-layer coverage ($96.32\%$ vs $97.90\%$), which we report and explain as
  a composition effect.

  \item \textbf{Geometric terminal takeover is a hybrid rescue, not RL learning.}
  Where used, early geometric takeover is a deployment-time rescue of the terminal
  docking failure (occupying $\sim\!27$--$43\%$ of steps / $\sim\!20\%$ of path for
  collapsed seeds); the RL policy itself is not improved. We describe it as a hybrid
  system, complementary to warm-start (which treats the full-speed-circling basin),
  not as ``the RL learned to dock.''

  \item \textbf{[PENDING]} multi-target empirical backing; a hard-dataset stress
  test to make the shield's value visible where the current benchmark is too easy;
  external CBF/VO Pareto baselines.
\end{itemize}

% =====================================================================
\section{Conclusion}
% [DRAFT]
We presented a projection-based, continuous-action COLREGs shield with an
honestly-scoped provable safety layer, filling the previously empty intersection of
continuous action, provable direction compliance, and COLREGs. We are explicit
about the boundary of the guarantee and about a training-stability limitation and
its warm-start remedy. Future work: multi-target empirical validation, a
multi-parent seed-variance study, yaw-axis smoothing, and external
(CBF/VO) Pareto baselines.

\section*{References}
% [PENDING: BibTeX; anchors: Krasowski \& Althoff 2024; Markgraf et al. 2026;
%  Meyer et al. 2020; Florensa et al. 2017; JSRL 2022; AWAC 2020; RLPD 2023;
%  Henderson et al. 2018; Agarwal et al. 2021.]

\end{document}
